In this section, we prove Theorem \ref{thm:dmft-approx}. We will omit arguments that are similar to ones in \cite{fan2025dynamicalI} for brevity, and focus on the differences in our setting of a non-Gaussian disorder matrix $\X$.

\subsection{Kernel spaces and kernel mappings}\label{subsec:spaces}

Fixing two constants $K_0,K_1>0$, we define ``envelope'' functions 
$\Phi_{R_\theta},\Phi_{R_g},\Phi_{C_\theta},\Phi_{C_g}$ on $[0,\infty)$ via
\begin{align}
\frac{\d}{\d t}\Phi_{R_\theta}(t) &= K_0\left(\Phi_{R_\theta}(t) + \int_0^t
\Phi_{R_g}(t-s)\Phi_{R_\theta}(s)\d s\right), \quad
\Phi_{R_\theta}(0)=K_0,\label{eq:envelope-R-theta}\\
\Phi_{R_g}(t) &= \sum_{p \geq 1} K_1^{p+1} \Phi_{R_\theta}^{\ast p}(t), \label{eq:envelope-R-g}\\
\frac{\d}{\d t}\Phi_{C_\theta}(t) &= K_0\left(\Phi_{C_\theta}(t)
+\int_0^t \Phi_{R_g}(t-s)[\Phi_{C_\theta}(s)+\Phi_{C_\theta}(t)]\d s
+ \Phi_{C_g}(t)+1\right), \quad
\Phi_{C_\theta}(0)=K_0,\label{eq:envelope-C-theta}\\
\sqrt{\Phi_{C_g}(t)} &= \sum_{p \geq 0} K_1^{p+1} \Phi_{R_\theta}^{\ast p} \ast
\sqrt{\Phi_{C_\theta}}(t),\label{eq:envelope-C-g}
\end{align}
where, for $a,b:[0,\infty) \to \R$, $[a*b](t)=\int_0^t a(t-s)b(s)\d s$.
It may be verified through a Laplace transform argument that 
this system has a unique solution, satisfying 
$\Phi_{R_\theta},\Phi_{R_g},\Phi_{C_\theta},\Phi_{C_g} \in S(\lambda)$ 
for all sufficiently large $\lambda>0$ where
\[
  S(\lambda)=
  \left\{
    \Phi:[0,\infty)\to \R,\;\Phi(t),\frac{\d}{\d t}\Phi(t) \geq 0 \text{ for all } t \geq 0,\;
    \int_0^\infty e^{-\lambda t}\Phi(t)\d t<\infty
  \right\}.
\]
We omit details of this verification, and refer to
\cite[Lemma 3.1]{fan2025dynamicalI} for an analogous argument.

We define from these envelope functions the following kernel spaces.

\begin{definition}\label{def:S}
Let $T>0$, let $D=\{d_1,\ldots,d_m\} \subset [0,T]$ be a finite subset,
and call \[[0,d_1),[d_1,d_2),\ldots,[d_{m-1},d_m),[d_m,T]\] 
the maximal intervals of $[0,T] \setminus D$. 

Fixing $T$-dependent
constants $K_0,K_1,K_2,K_3>0$ where $K_0,K_1$ define the above envelopes,
let $\mathcal{S}_\theta^+(T,D)$ be the space of all
pairs $(C_\theta,R_\theta)$ such that:
    \begin{enumerate}
    \item $C_\theta:[0,T] \times [0,T] \to \R$ is a positive-semidefinite
covariance kernel satisfying
\[C_\theta(t,t) \leq \Phi_{C_\theta}(t) \text{ for all } 0\leq t\leq T,\]
and for any maximal interval $I$ of $[0,T]\setminus D$,
        \begin{equation}\label{eq:Ctheta-cont}
            |C_\theta(t,t)+C_\theta(s,s)-2C_\theta(t,s)| \leq  K_2|t-s|
\text{ for all } t,s \in I.
        \end{equation}
        \item $R_\theta:[0,T] \times [0,T] \to \R$ is a function satisfying
\[|R_\theta(t,s)| \leq \Phi_{R_\theta}(t-s) \text{ for all } 0 \leq s \leq t
\leq T, \qquad R_\theta(t,s)=0 \text{ for all } s>t,\]
and for any two maximal intervals $I,I'$ of $[0,T]\setminus D$,
\begin{equation}\label{eq:Rtheta-cont}
|R_\theta(t,s)-R_\theta(t',s')| \leq K_2(|t-t'|+|s-s'|) \text{ for all } t,t'\in
I,~s,s'\in I' \text{ with } s\le t \text{ and } s'\le t'.
\end{equation}
\end{enumerate}

Let $\mathcal{S}_g(T,D)$ be the space of all pairs
$(C_g,R_g)$ such that:
\begin{enumerate}
    \item $C_g:[0,T] \times [0,T] \to \R$ is a symmetric function (not necessarily
positive-semidefinite) satisfying
\[C_g(t,t) \leq \Phi_{C_g}(t) \text{ for all } 0\leq t\leq T,\]
and for any two maximal intervals $I,I'$ of $[0,T]\setminus D$,
        \begin{equation}\label{eq:Cg-cont1}
|C_g(t,t)+C_g(s,s)-2C_g(t,s)| \leq K_2K_3|t-s|
\text{ for all } t,s \in I,
\end{equation}
\begin{equation}\label{eq:Cg-cont2}
            |C_g(t,s)-C_g(t',s')| \leq K_2K_3(\sqrt{|t-t'|}+\sqrt{|s-s'|})
\text{ for all } t,t' \in I,~s,s' \in I'.
        \end{equation}
      \item $R_g:[0,T] \times [0,T] \to \R$ is a function satisfying
\[|R_g(t,s)| \leq \Phi_{R_g}(t-s) \text{ for all } 0 \leq s \leq t \leq T,
\qquad R_g(t,s)=0 \text{ for all } s>t,\]
and for any two maximal intervals $I,I'$ of $[0,T]\setminus D$,
\begin{equation}\label{eq:Rg-cont}
|R_g(t,s)-R_g(t',s')| \leq K_2K_3(|t-t'|+|s-s'|) \text{ for all } t,t'\in
I,~s,s'\in I' \text{ with } s\le t \text{ and } s'\le t'.
\end{equation}
\end{enumerate}
We denote by $\cS_g^+(T,D) \subset \cS_g(T,D)$
the subset of pairs $(C_g,R_g)$ where $C_g$ is positive-semidefinite.
\end{definition}

These spaces depend on the constants $K_0,K_1,K_2,K_3>0$, although we make this
dependence implicit in the notation. The final spaces $\cS_\theta^+,\cS_g^+$ in
Theorem \ref{thm:dmft-approx}(a) will be the spaces of kernels
$C_\theta,R_\theta,C_g,R_g$ on $[0,\infty) \times [0,\infty)$ whose restrictions
to $[0,T] \times [0,T]$
belong to $\cS_\theta^+(T,\emptyset)$ and $\cS_g^+(T,\emptyset)$ with
$D=\emptyset$ for every $T>0$.

We define a mapping
\[\cT_{\theta \to g}:\cS_\theta^+(T,D) \to \cS_g(T,D),
\qquad (C_\theta,R_\theta) \mapsto (C_g,R_g)\]
by (\ref{eq:dmft-Cg}--\ref{eq:dmft-Rg}):
\[C_g=\sum_{p,q \geq 0} \kappa_{p+q+2}\,
R_\theta^{\ast p} * C_\theta * \bar R_\theta^{\ast q},
\qquad
R_g=\sum_{p \geq 1} \kappa_{p+1}\,R_\theta^{\ast p}.\]
Note that $C_g$ is not necessarily positive-semidefinite.
Restricting to positive-semidefinite $C_g$, we define a reverse mapping
\[\cT_{g \to \theta}:\cS_g^+(T,D) \to \cS_\theta^+(T,\emptyset),
\qquad (C_g,R_g) \mapsto (C_\theta,R_\theta)\]
by (\ref{eq:dmft-Ctheta}--\ref{eq:dmft-Rtheta}):
\[C_\theta(t,s)=\E[\theta^t \theta^s],
\qquad R_\theta(t,s)=\E\left[\frac{\d\theta^t}{\d g^s}\right]\]
where $\{\theta^t\}_{t \geq 0}$ and $\{\frac{\partial \theta^t}{\partial
g^s}\}_{t \geq s \geq 0}$ solve (\ref{eq:dmft-theta}--\ref{eq:dmft-response})
with $R_g$ and $\{g^t\}_{t \geq 0} \sim \GP(0,C_g)$. The following lemma
verifies that these mappings are well-defined.

\begin{lemma}\label{lemma:well_defined}
Fix any $T>0$. Then for all sufficiently large ($T$-dependent)
constants $K_0,K_1,K_2,K_3>0$ defining the spaces of Definition \ref{def:S},
the following holds:

For any finite subset $D \subset [0,T]$, given any $(C_g,R_g) \in
\cS_g^+(T,D)$, the system (\ref{eq:dmft-theta}--\ref{eq:dmft-response}) has a
unique solution, and $\cT_{g\to \theta}(C_g,R_g) \in \cS_\theta^+(T,\emptyset)$.
Given any $(C_\theta,R_\theta) \in \cS_\theta^+(T,D)$, 
the series (\ref{eq:dmft-Cg}--\ref{eq:dmft-Rg}) are pointwise convergent, and
$\cT_{\theta\to g}(C_\theta,R_\theta) \in \cS_g(T,D)$.
\end{lemma}
\begin{proof}
Given $(C_g,R_g) \in \cS_g^+(T,D)$, existence and uniqueness of $\{\theta^t\}_{t
\in [0,T]}$ and $\{\frac{\partial \theta^t}{\partial g^s}\}_{0 \leq s \leq t
\leq T}$ solving (\ref{eq:dmft-theta}--\ref{eq:dmft-response}) follows from the
same argument as \cite[Lemma~3.3]{fan2025dynamicalI}, which we omit here for brevity.
Therefore, $\cT_{g\to \theta}$ is well-defined.

Let $(C_\theta,R_\theta)=\cT_{g \to \theta}(C_g,R_g)$.
We check that $(C_\theta,R_\theta) \in \cS_\theta^+(T,\emptyset)$:
First, according to \eqref{eq:dmft-response},
\begin{align*}
 & \left| \frac{\d }{\d t} \frac{\partial \theta^t}{\partial g^s}\right|
 \leq \|f\|_\Lip \left| \frac{\partial \theta^t}{\partial g^s}\right| + \int_s^t
\Phi_{R_g}(t-r) \left|\frac{\partial\theta^r}{\partial g^s}\right|\d r
\end{align*} 
where $\|f\|_\Lip<\infty$ is the Lipschitz constant of $f$.
Noting that $\frac{\partial \theta^s}{\partial g^s}=1$ and
comparing with \eqref{eq:envelope-R-theta}, this implies for $K_0>0$ large enough
that
\begin{align}\label{eq:bound-response-process} 
\left|\frac{\partial \theta^t}{\partial g^s}\right| \leq \Phi_{R_\theta}(t-s).
\end{align}
Then by definition of $R_\theta$ in \eqref{eq:dmft-Rtheta},
$|R_\theta(t,s)| \le \Phi_{R_\theta}(t-s)$ for all $0\le s\le t\le T$.
Similarly, for any $0 \le s,t,t' \le T$ with $s \leq t' \leq t$, we have
\begin{align}
    |R_\theta(t,s) - R_\theta(t',s)| &\le \E\left|\frac{\partial
\theta^t}{\partial g^s} - \frac{\partial \theta^{t'}}{\partial g^s}\right| \le
\E \int_{t'}^t \left| \frac{\d }{\d r} \frac{\partial \theta^r}{\partial g^s}
\right|\d r\notag\\
    &\le \E \int_{t'}^t \left(\|f\|_\Lip \left| \frac{\partial
\theta^r}{\partial g^s}\right| +  \int_s^r \Phi_{R_g}(r-r')
\left|\frac{\partial\theta^{r'}}{\partial g^s}\right|\d r'\right)\d r
    \le K_2|t-t'|,\label{eq:Rthetacontt}
\end{align}
the last inequality holding by \eqref{eq:bound-response-process} for some
large enough ($T$-dependent) $K_2>0$.
For continuity in \(s\), by the response equation \eqref{eq:dmft-response}, for
\(t \ge s\),
\[
\frac{\partial \theta^t}{\partial g^s}
=
1+
\int_s^t
\left[
f'(\theta^u)\frac{\partial \theta^u}{\partial g^s}
+
\int_s^u R_g(u,r)\frac{\partial \theta^r}{\partial g^s}\,\d r
\right]\d u.
\]
Consider any $t \geq s' \geq s$. Then
\[
\begin{aligned}
\left|
\frac{\partial \theta^{s'}}{\partial g^s}-1
\right|
&\le
\int_s^{s'}
\|f\|_\Lip \left|\frac{\partial \theta^u}{\partial g^s}\right|\,\d u
+
\int_s^{s'}\int_s^u |R_g(u,r)|
\left|\frac{\partial \theta^r}{\partial g^s}\right|
\,\d r\,\d u  \le C |s'-s|,
\end{aligned}
\]
the last holding by
\eqref{eq:bound-response-process} and $|R_g(u,r)| \le \Phi_{R_g}(T)$
for some large enough ($T$-dependent) constant $C>0$.
Then subtracting the two response equations for
\(\frac{\partial \theta^t}{\partial g^s}\) and
\(\frac{\partial \theta^t}{\partial g^{s'}}\) gives
\begin{align*}
&\left|
\frac{\partial \theta^t}{\partial g^s}
-\frac{\partial \theta^t}{\partial g^{s'}}
\right|
\le
\left|
\frac{\partial \theta^{s'}}{\partial g^s}
-1 \right| +
\int_{s'}^t
\|f\|_\Lip
\left|
\frac{\partial \theta^u}{\partial g^s}
-
\frac{\partial \theta^u}{\partial g^{s'}}
\right|
\,\d u \\
&\quad+
\int_{s'}^t\int_{s'}^u
|R_g(u,r)|
\left|
\frac{\partial \theta^r}{\partial g^s}
-
\frac{\partial \theta^r}{\partial g^{s'}}
\right|
\,\d r\,\d u +
\int_{s'}^t\int_s^{s'}
|R_g(u,r)|
\left|
\frac{\partial \theta^r}{\partial g^s}
\right|
\,\d r\,\d u.
\end{align*}
The first and last terms are bounded by \(C|s'-s|\) for some $C>0$, and the
second and third terms by $C \sup_{u \in [s',t]} |
\frac{\partial \theta^u}{\partial g^s}
-\frac{\partial \theta^u}{\partial g^{s'}}|$.
Then by Gr\"onwall's inequality,
\[
\left|
\frac{\partial \theta^t}{\partial g^s}
-
\frac{\partial \theta^t}{\partial g^{s'}}
\right|
\le
K_2|s'-s|
\]
for some sufficiently large $K_2>0$. Taking expectations gives
\begin{equation}\label{eq:Rthetaconts}
\left|R_\theta(t,s)-R_\theta(t,s')\right| \le K_2|s'-s|.
\end{equation}
Combining \eqref{eq:Rthetacontt} and \eqref{eq:Rthetaconts} shows the continuity property \eqref{eq:Rtheta-cont}.
This verifies all properties for $R_\theta$ in Definition \ref{def:S}.

For the covariance kernel $C_\theta$, applying It\^o's formula to 
the equation \eqref{eq:dmft-theta} gives
  \begin{align*}
      \frac{\d}{\d t} \E (\theta^t)^2 &= 2\E[\theta^t f(\theta^t)] + 2\int_0^t
R_g(t,s)\E[\theta^t\theta^s]\d s + 2\E[\theta^t g^t] + 2\gamma\\
      &\le K_0\left(\E(\theta^t)^2+\int_0^t \Phi_{R_g}(t-s) [\E (\theta^s)^2
+ \E (\theta^t)^2] \d s + C_g(t,t)+1\right),
  \end{align*}
the second inequality holding for large enough $K_0>0$
since $\{g^t\}_{t \geq 0} \sim \GP(0,C_g)$, $f$ is Lipschitz, and $|R_g(s,r)|\le
\Phi_{R_g}(s-r)$.
Comparing with \eqref{eq:envelope-C-theta}, this implies 
$C_\theta(t,t)=\E(\theta^t)^2 \le \Phi_{C_\theta}(t)$.
Similarly, for all $s,s',t,t' \in [0,T]$, one can show
\[
|C_\theta(t,t)+C_\theta(s,s)-2C_\theta(t,s)| = \E (\theta^t - \theta^s)^2
\le K_2|t-s|
\]
for large enough $K_2>0$. This verifies all properties for $C_\theta$ in
Definition \ref{def:S} with $D=\emptyset$, so $(C_\theta,R_\theta) \in \cS_\theta^+(T,\emptyset)$.

For the map $\cT_{\theta\to g}$, given $(C_\theta,R_\theta)\in
\cS_\theta^+(T,D)$, note that Proposition \ref{prop:freecumulants}
implies $|\kappa_p| \leq K_1^p$ for any sufficiently large $K_1>0$ and all
$p \geq 1$, and $|R_\theta(t,s)| \leq \Phi_{R_\theta}(t-s)$. Then
\[\left|\sum_{p \geq 1} \kappa_{p+1} R_\theta^{*p}(t,s)\right|
\leq \sum_{p \geq 1} K_1^{p+1} \Phi_{R_\theta}^{*p}(t-s)
\leq \sum_{p \geq 1} K_1^{p+1}\Phi_{R_\theta}(T)^p
\frac{(t-s)^{p-1}}{(p-1)!}<\infty,\]
so the series \eqref{eq:dmft-Rg} is pointwise convergent. Similarly
\eqref{eq:dmft-Cg} is pointwise convergent, so
$\cT_{\theta\to g}$ is well-defined.

Let $(C_g,R_g)=\cT_{\theta\to g}(C_\theta,R_\theta)$. 
Applying $|\kappa_p| \leq K_1^p$, $|R_\theta(t,s)| \leq \Phi_{R_\theta}(t-s)$,
and
\[
|C_\theta(t,s)| \le \sqrt{C_\theta(t,t)}\sqrt{C_\theta(s,s)} \le \sqrt{\Phi_{C_\theta}(t)}\sqrt{\Phi_{C_\theta}(s)},
\]
we can verify that
\begin{align*}
    |R_g(t,s)| &\le \sum_{p \geq 1} K_1^{p+1}\,\Phi_{R_\theta}^{\ast p}(t-s) =
\Phi_{R_g}(t-s),\\
    |C_g(t,t)| &\le \sum_{p,q \geq 0} K_1^{p+q+2}\,
(\Phi_{R_\theta}^{\ast p} \ast \sqrt{\Phi_{C_\theta}})(t) (\Phi_{R_\theta}^{\ast
q} \ast \sqrt{\Phi_{C_\theta}})(t)\\
&=\left(\sum_{p\ge 0} K_1^{p+1}\,
(\Phi_{R_\theta}^{\ast p} \ast \sqrt{\Phi_{C_\theta}})(t)\right)^2 = \Phi_{C_g}(t).
\end{align*}

For any maximal interval $I$ of $[0,T] \setminus D$ and any $s,t \in I$
with $t \geq s$, note that
\begin{align}
&|C_g(t,t)+C_g(s,s)-2C_g(t,s)|\notag\\
&=\Bigg|\kappa_2\Big(C_\theta(t,t)+C_\theta(s,s)-2C_\theta(t,s)\Big)
+\sum_{p\ge 1}
  \kappa_{p+2}
  \int_0^t (R_\theta^{\ast p}(t,r)-R_\theta^{\ast p}(s,r))(C_\theta(r,t)
-C_\theta(r,s))\d r \notag\\
&\hspace{0.5in}+ \sum_{q\ge 1}
  \kappa_{q+2}
  \int_0^t \d u\,
  (C_\theta(t,u)-C_\theta(s,u))
(R_\theta^{\ast q}(t,u)-R_\theta^{\ast q}(s,u))\notag\\
&\hspace{0.5in}+\sum_{p,q\ge 1}
  \kappa_{p+q+2}
  \int_0^t \d r \int_0^t \d u\,
    (R_\theta^{\ast p}(t,r)-R_\theta^{\ast p}(s,r))C_\theta(r,u)
    (R_\theta^{\ast q}(t,u)-R_\theta^{\ast q}(s,u))\Bigg|.\label{eq:Cgdiffexpr}
\end{align}
For any \(p\ge 1\) and $t \ge s\ge r$,
\[
R_\theta^{\ast p}(t,r)-R_\theta^{\ast p}(s,r)
=
\int_r^s
\bigl(R_\theta(t,u)-R_\theta(s,u)\bigr)
R_\theta^{\ast(p-1)}(u,r)\,\d u+
\int_s^t
R_\theta(t,u)R_\theta^{\ast(p-1)}(u,r)\,\d u.
\]
Therefore, combined with 
\[
|R_\theta(t,u) - R_\theta(s,u)| \leq K_2 |t-s|, \qquad |R_\theta^{\ast p}(t,s)| \le \Phi_{R_\theta}^{\ast p}(t-s) \le \frac{\Phi_{R_\theta}(T)^p (t-s)^{p-1}}{(p-1)!},
\]
we have that
\begin{align*}
|R_\theta^{\ast p}(t,r)-R_\theta^{\ast p}(s,r)|
&\le \int_r^s
K_2|t-s|
\Phi_{R_\theta}^{\ast (p-1)}(u-r)\,\d u 
+
\int_s^t \Phi_{R_\theta}(t-u)\Phi_{R_\theta}^{\ast (p-1)}(u-r)
 \d u\\
 &\le (K_2T+\Phi_{R_\theta}(T)) \frac{\Phi_{R_\theta}(T)^{p-1}T^{p-1}}{(p-1)!} |t-s|. 
\end{align*}
This bound holds also for $r \in [s,t]$, noting in this case that
\[R_\theta^{\ast p}(t,r)-R_\theta^{\ast p}(s,r)
=R_\theta^{\ast p}(t,r)=\int_r^t R_\theta(t,u)R_\theta^{*(p-1)}(u,r)\,\d u\]
and applying the same argument. Applying this bound for summands of
\eqref{eq:Cgdiffexpr} where either $p \geq 1$ or $q \geq 1$, and applying
$|C_\theta(t,t)+C_\theta(s,s)-2C_\theta(t,s)| \leq K_2|t-s|$
for $p=q=0$, we obtain
\[|C_g(t,t)+C_g(s,s)-2C_g(t,s)| \leq K_2K_3|t-s|\]
for all large enough $K_2,K_3>0$, which is the continuity property
\eqref{eq:Cg-cont1}.

Since $C_\theta$ is positive-semidefinite, writing $u \sim \GP(0,C_\theta)$,
for any two maximal intervals $I,I'$ of $[0,T] \setminus D$ and any
$t,t' \in I$ and $s,s' \in I'$, we have also
\[|C_\theta(t,s)-C_\theta(t',s')|
\leq |\E[(u^t-u^{t'})u^s]|+|\E[u^{t'}(u^s-u^{s'})]|
\leq \sqrt{\Phi_{C_\theta}(T)} \cdot K_2(\sqrt{|t-t'|}+\sqrt{|s-s'|}),\]
by Cauchy-Schwarz. Then applying a similar argument as above shows
\[|C_g(t,s)-C_g(t',s')| \leq K_2K_3(\sqrt{|t-t'|}+\sqrt{|s-s'|})\]
for all large enough $K_2,K_3>0$,
which is the continuity property \eqref{eq:Cg-cont2}.
Finally, for any maximal intervals $I,I'$ of
$[0,T] \setminus D$ and $s,s' \in I$, $t,t' \in I'$ with $s \leq t$ and $s' \leq
t'$, a similar argument as above using $|R_\theta(t,s)-R_\theta(t',s')| \leq
K_2(|t-t'|+|s-s'|)$ checks that
\[|R_g(t,s)-R_g(t',s')| \leq K_2K_3(|t-t'|+|s-s'|)\]
for large enough $K_2,K_3>0$, which is \eqref{eq:Rg-cont}.
Thus $(C_g,R_g)\in \cS_g(T,D)$.
\end{proof}

Henceforth, given $T>0$, we will fix $T$-dependent constants
$K_0,K_1,K_2,K_3>0$ large enough such that Lemma \ref{lemma:well_defined}
holds.

\subsection{Contractivity and uniqueness of fixed point with projection}

Next, we define a metric on $\cS_\theta^+$ and $\cS_g$. Fix $\lambda>0$.
For a pair of response kernels $(R_1,R_2) \equiv (R_{\theta,1},R_{\theta,2})$ 
or $(R_1,R_2) \equiv (R_{g,1},R_{g,2})$ satisfying the conditions of
Definition \ref{def:S}, we define
\begin{align}\label{eq:metric-response}
  d_\lambda(R_1, R_2) = \sup_{0 \leq s \leq t \leq T} e^{-\lambda
t}|R_1(t,s)-R_2(t,s)|.
\end{align} 
For a pair of (not necessarily positive-semidefinite) correlation kernels
$(C_1,C_2) \equiv (C_{\theta,1},C_{\theta,2})$ 
or $(C_1,C_2) \equiv (C_{g,1},C_{g,2})$ satisfying the conditions of
$\cS_\theta^+$/$\cS_g$ in Definition \ref{def:S}, we define
\begin{align}\label{eq:hilbertinnerproduct}
\langle C_1, C_2 \rangle_{\lambda}=\int_0^T \int_0^T e^{-2\lambda
(t\vee s)}C_1(t,s) C_2(t,s)\,\d t\,\d s,
\qquad \|C\|_\lambda=\langle C,C \rangle_{\lambda}^{1/2}
 \end{align} 
and the metric
\begin{align}
\label{eq:metric-covariance}
  d_\lambda(C_1, C_2)=\|C_1-C_2\|_\lambda= \sqrt{\int_0^T \int_0^T e^{-2\lambda (t\vee s)}(C_1(t,s) - C_2(t,s))^2 \,\d t \,\d s}.
\end{align}
Finally, for either
$(C,R) \equiv (C_\theta,R_\theta)$ or $(C,R) \equiv (C_g,R_g)$, we set
\[d_\lambda((C_1,R_1),(C_2,R_2))
=d_\lambda(R_1,R_2)+d_\lambda(C_1,C_2).\]
By the uniform continuity properties of the kernels on $[0,T] \setminus D$
in Definition \ref{def:S},
this defines a valid metric $d_\lambda$ on both $\cS_\theta^+$ and
$\cS_g$.

\begin{lemma}[Contractivity of $\cT_{g \to \theta}$]\label{lem:contraction-Rtheta}
Let $X_i=(C_{g,i},R_{g,i})\in \mathcal{S}_{g}^+(T,D)$ for $i=1,2$, and let
$Y_i=(C_{\theta,i},R_{\theta,i})=\cT_{g \to \theta}(C_{g,i},R_{g,i})$.
Then for any fixed $0<\alpha<1$ and all sufficiently large $\lambda>0$
(depending on $\alpha,T$),
\begin{align}\label{equation:contraction_Rfix}
d_\lambda(Y_1,Y_2) \leq \alpha\cdot d_\lambda(X_1,X_2).
\end{align}
\end{lemma}
\begin{proof}
We first consider $C_{g,1}=C_{g,2} \equiv C_g$ and different $R_{g,1},R_{g,2}$.
Coupling the processes $\{\theta_i^t\}_{t \geq 0}$ 
and $\{\frac{\partial \theta_i^t}{\partial g^s}\}_{t \geq s \geq 0}$ 
solving (\ref{eq:dmft-theta}--\ref{eq:dmft-response}) with $R_{g,i}$
by the same realizations of $\{g^t\}_{t \in [0,T]} \sim \GP(0,C_g)$ and
$\{b^t\}_{t \in [0,T]}$, a Gr\"onwall
argument (c.f.\ \cite[Lemma 3.7]{fan2025dynamicalI}) shows,
for all sufficiently large $\lambda>0$ and a constant $C>0$ not depending on
$\lambda$,
\begin{align*}
\sup_{0 \leq t \leq T} e^{-2\lambda t}\E(\theta_1^t-\theta_2^t)^2 \leq
\frac{C}{\lambda}\,d_\lambda(R_{g,1},R_{g,2})^2,\\
\sup_{0 \leq s \leq t \leq T}
e^{-\lambda t}
\E \left|\frac{\partial \theta_1^t}{\partial g^s} - \frac{\partial
\theta_2^t}{\partial g^s}\right| \leq \frac{C}{\lambda}
\,d_\lambda(R_{g,1},R_{g,2}).
\end{align*}
The second inequality implies
$d_\lambda(R_{\theta,1},R_{\theta,2}) \leq \frac{C}{\lambda}
d_\lambda (R_{g,1},R_{g,2})$. By Cauchy-Schwarz,
\begin{align*}
|C_{\theta,1}(t,s)-C_{\theta,2}(t,s)|^2 &\leq 2\E(\theta_1^t-\theta_2^t)^2\E(\theta_1^s)^2 + 2\E(\theta_2^t)^2\E(\theta_1^s-\theta_2^s)^2\\
&\leq 2\Phi_{C_\theta}(T)\left( \E(\theta_1^t-\theta_2^t)^2 + \E(\theta_1^s-\theta_2^s)^2 \right). \end{align*}
Integrating both sides, we have
\[d_{\lambda}(C_{\theta,1}, C_{\theta,2})^2
=\int_0^T\int_0^T e^{-2\lambda(t \vee s)}
|C_{\theta,1}(t,s)-C_{\theta,2}(t,s)|^2\,\d t\, \d s
\leq 4\Phi_{C_\theta}(T)T^2
\sup_{0 \leq t \leq T} e^{-2\lambda t}\E(\theta_1^t-\theta_2^t)^2,\]
so the first inequality implies 
$d_{\lambda}(C_{\theta,1},C_{\theta,2}) \leq \frac{C'}{\sqrt{\lambda}}
d_\lambda (R_{g, 1}, R_{g, 2})$. Thus, for any $\alpha \in (0,1)$ and all large
enough $\lambda>0$,
\begin{equation}\label{eq:contractionsameCg}
d_\lambda(Y_1,Y_2) \leq \alpha \cdot d_\lambda(R_{g,1},R_{g,2}).
\end{equation}

We next consider $R_{g,1}=R_{g,2} \equiv R_g$ and different $C_{g,1},C_{g,2}$.
For $u \in [0,1]$, define the interpolation $C_u=(1-u)C_{g,1}+uC_{g,2}$,
and let $g_u \sim \GP(0,C_u)$.
Note that $C_u$ induces a positive-semidefinite, trace-class covariance operator
$Q_u$ on $\mathcal{H}:=L^2([0,T])$, given by
\[
    (Q_u\varphi)(r)
    :=
    \int_0^T C_u(r,r')\varphi(r')\,\d r',
    \qquad
    \varphi\in \mathcal{H}.
\]
Let $\{e_k\}_{k\ge 1}$ be an orthonormal basis of
$\mathcal{H}$, and let $\{\xi_k\}_{k\ge 1}$ be i.i.d.~standard Gaussian random
variables. Then we may realize $g_u \sim \GP(0,C_u)$ as
$g_u=\sum_{k\ge 1}Q_u^{1/2}e_k \cdot \xi_k$.

Fixing an integer $N \geq 1$, define $\Pi_N : {\mathcal H}\to {\sp(e_1,\ldots, e_N)}$ as the orthogonal projection onto the subspace spanned by $e_1,\ldots,e_N$. Consider the truncation
$g_{u,N}:=\Pi_N g_u$, 
and define $\theta_{g_{u,N}}$ by
\[
    \d \theta_{g_{u,N}}^t
    =
        \left[f(\theta_{g_{u,N}}^t)
        +
        \int_0^t R_g(t,s)\theta_{g_{u,N}}^s \d s
        +
        g_{u,N}^t \right]
    \d t + \sqrt{2\gamma}\,\d b^t, \qquad \theta_{g_{u,N}}^0=\theta^0.
\] 
Let us compute
$\frac{\d}{\d u} \E[\theta_{g_{u,N}}^t\theta_{g_{u,N}}^s]$:
For any $h\in \mathcal{H}$, define $\{\theta_h^t\}_{t \geq 0}$ by
\[
    \theta_{h}^t
    =
    \theta^0
    +
    \int_0^t
    \left[
        f(\theta_{h}^s)
        +
        \int_0^s R_g(s,r)\theta_{h}^r \d r
        +
        {h}^s
    \right]
    \d s + \sqrt{2\gamma}\,b^t.
\]
By \cite[Theorem 2.5]{wang2017stochastic}, for any $h,k,l \in \cH$,
the above equation has a unique solution $\{\theta_h^t\}_{t \geq 0}$, and
there exist directional derivative processes
$D_k\theta_h$ and $D_{kl}^2 \theta_h$ on $[0,T]$
for which
\[\lim_{\eps \to 0}
\E_{\{b^t\}}\left(\frac{\theta_{h+\eps k}^t-\theta_h^t}{\eps}
-D_k\theta_h^t\right)^2=0,
\qquad \lim_{\eps \to 0}
\E_{\{b^t\}}\left(\frac{D_k\theta_{h+\eps l}^t-D_k\theta_h^t}{\eps}
-D_{kl}^2\theta_h^t\right)^2=0\]
where $\E_{\{b^t\}}$ is the expectation over the Brownian motion
$\{b^t\}_{t \geq 0}$. Then 
$\E_{\{b^t\}}[\theta_{h+\eps k+\eps'l}^t
\theta_{h+\eps k+\eps'l}^s]$ is twice continuously-differentiable in
$(\eps,\eps')$, and 
\begin{equation}\label{eq:thetatssecondder}
\partial_\eps\partial_{\eps'}
\E_{\{b^t\}}[\theta_{h+\eps k+\eps'l}^t
\theta_{h+\eps k+\eps'l}^s]\big|_{\eps=\eps'=0}
=\E_{\{b^t\}}\Big[(D_{kl}^2\theta_h^t)\theta_h^s
+\theta_h^t (D_{kl}^2\theta_h^s)+(D_k\theta_h^t)(D_l\theta_h^s)
+(D_l\theta_h^t)(D_k\theta_h^s)\Big].
\end{equation}
Furthermore by \cite[Theorem 2.5]{wang2017stochastic}, $D_k\theta_h$ and
$D_{kl}^2\theta_h$ are the unique solutions of
\begin{align*}
    D_k\theta_h^t
    &=
    \int_0^t
    \left[
        f'(\theta_h^s)D_k\theta_h^s
        +
        \int_0^s R_g(s,r) D_k\theta_h^r \d r
        +
        k^s
    \right]
    \d s,\\
    D_{kl}^2\theta^t
    &=
    \int_0^t
    \left[
        f''(\theta_h^s)
        (D_k\theta_h^s)(D_l\theta_h^s)
        +f'(\theta_h^s)D_{kl}^2 \theta_h^s
        + \int_0^s R_g(s,r) D_{kl}^2\theta_h^r \d r
    \right]
    \d s.
\end{align*}
Equivalently, define first-order and second-order response kernels $J_h,K_h$ by
\begin{align*}
    J_h(t,r)
    &=
    1
    +
    \int_r^t
    \left[
        f'(\theta_h^s)J_h(s,r)
        +
        \int_r^s R_g(s,r')J_h(r',r)\d r
    \right]
    \d s\\
    K_h(t,r,r')
    &=
    \int_{r \vee r'}^t
    \bigg[
        f''(\theta_h^s)J_h(s,r)J_h(s,r')
        +
        f'(\theta_h^s)K_h(s,r,r')
        +
        \int_{r \vee r'}^s R_g(s,s') K_h(s',r,r') \d s' 
    \bigg]
    \d s
\end{align*}
for $t \geq \max(r,r')$, and $J_h(t,r)=0$ and $K_h(t,r,r')=0$ if $t<r$ and/or
$t<r'$. Then it may be checked that the above solutions are given by
\begin{equation}\label{eq:second-deriv-representation}
D_k\theta_h^t=\int_0^T J_h(t,r) k^r\,\d r,
\qquad
D_{kl}^2\theta_h^t=\int_0^T \int_0^T K_h(t,r,r')k^r l^{r'}\,\d r\,\d r'.
\end{equation}
Hence \eqref{eq:thetatssecondder} takes the form
\begin{align*}
\partial_\eps\partial_{\eps'}
\E_{\{b^t\}}[\theta_{h+\eps k+\eps'l}^t
\theta_{h+\eps k+\eps'l}^s]\big|_{\eps=\eps'=0}
&=\int_0^T \int_0^T H_h(t,s;r,r')k^r l^{r'}\,\d r\,\d r',
\end{align*}
where
\begin{align*}
    H_h(t,s;r,r')
    &:=
    K_h(t,r,r')\theta_h^s
    +
    K_h(s,r,r')\theta_h^t
    +
    J_h(t,r)J_h(s,r')
    +
    J_h(t,r')J_h(s,r).
\end{align*}
Then, by the finite-dimensional Gaussian interpolation formula
(see e.g.\ \cite[Lemma 7.2.7]{vershynin2018high}),
\begin{align}
    \frac{\d}{\d u}\E[\theta_{g_{u,N}}^t\theta_{g_{u,N}}^s]
&=\frac{\d}{\d u}\E[\E_{\{b^t\}}[\theta_{g_{u,N}}^t\theta_{g_{u,N}}^s]]\notag\\
    &=
    \frac12
    \sum_{i,j=1}^N
    \bigl\langle (Q_1-Q_0)e_i,e_j\bigr\rangle_{\mathcal{H}}
    \,
    \E
\left[\partial_\eps\partial_{\eps'}
\E_{\{b^t\}}[\theta_{g_{u,N}+\eps e_i+\eps'e_j}^t
\theta_{g_{u,N}+\eps e_i+\eps'e_j}^s]\big|_{\eps=\eps'=0}\right]
    \notag\\
    &=
    \frac12
    \sum_{i,j=1}^N
    \bigl\langle (Q_1-Q_0)e_i,e_j\bigr\rangle_{\mathcal{H}}
    \,
    \E\left[
        \int_0^T\int_0^T
        H_{g_{u,N}}(t,s;r,r')
        e_i^re_j^{r'}
        \,\d r\,\d r'
    \right].\label{eq:duCg}
\end{align}

By the Parseval identity, as $N \to \infty$,
\[\sum_{i,j=1}^N
    \bigl\langle (Q_1-Q_0)e_i,e_j\bigr\rangle_{\mathcal{H}}
    e_i^re_j^{r'} \to (C_{g,2}-C_{g,1})(r,r')
\]
in $L^2([0,T] \times [0,T])$.
Also as $N \to \infty$, $g_{u,N} \to g_u$ in $L^2([0,T])$. Then
applying a Gr\"onwall argument to the definitions of the above processes
shows that $\theta_{g_{u,N}} \to \theta_{g_u}$ and
$J_{g_{u,N}}(t,\cdot) \to J_{g_u}(t,\cdot)$
in $L^2([0,T])$, and $K_{g_{u,N}}(t,\cdot,\cdot) \to
K_{g_u}(t,\cdot,\cdot)$ and
$H_{g_{u,N}}(t,s;\cdot,\cdot) \to H_{g_u}(t,s;\cdot,\cdot)$
in $L^2([0,T] \times [0,T])$. Then,
integrating \eqref{eq:duCg} over $u \in [0,1]$ and taking the limit $N \to
\infty$,
\[(C_{\theta,2}-C_{\theta,1})(t,s)
=\E[\theta_{g_1}^t\theta_{g_1}^s]
-\E[\theta_{g_0}^t\theta_{g_0}^s]
    =\int_0^1 \d u\left(
    \frac12
    \int_0^T\int_0^T
    (C_{g,2}-C_{g,1})(r,r')
    \E\left[
        H_{g_u}(t,s;r,r')
    \right]
    \,\d r\,\d r'\right).
\]
As $f',f''$ are bounded and $R_g$ is bounded on $[0,T]$,
a Gr\"onwall argument checks that $J_h(t,r)$ and $K_h(t,r,r')$ are both uniformly
bounded over $t,r,r' \in [0,T]$ and $h =g_u$ for all $u \in [0,1]$. Then
$\E[H_{g_u}(t,s;r,r')]$ is also uniformly bounded over $t,s,r,r' \in [0,T]$
and $u \in [0,1]$.  Furthermore, $\E[H_{g_u}(t,s;r,r')]=0$ unless $r,r' \leq
\max(t,s)$. Thus, there exists a $T$-dependent constant $C>0$ for which
\begin{align*}
|(C_{\theta,2}-C_{\theta,1})(t,s)|
&\leq C\int_0^{t \vee s}\int_0^{t \vee s}
    |(C_{g,2}-C_{g,1})(r,r')|\,\d r\,\d r'
\end{align*}
Then by Cauchy-Schwarz, for $T$-dependent constants $C',C''>0$ (not depending on
$\lambda$),
\begin{align*}
&d_\lambda(C_{\theta,1},C_{\theta,2})^2
=\int_0^T \int_0^T e^{-2\lambda(t \vee s)}
(C_{\theta,2}(t,s)-C_{\theta,1}(t,s))^2\d t\,\d s\\
&\leq C'\int_0^T\d t\int_0^T \d s\,
e^{-2\lambda(t \vee s)}
\int_0^{t \vee s}\d r\int_0^{t \vee s} \d r'
    (C_{g,2}(r,r')-C_{g,1}(r,r'))^2\\
&=C'\int_0^T \d r\int_0^T \d r'
e^{-2\lambda(r \vee r')}
(C_{g,2}(r,r')-C_{g,1}(r,r'))^2
\mathop{\int_0^T \d t\int_0^T\d s} \1_{t \vee s \geq r \vee r'}
e^{-2\lambda(t \vee s)+2\lambda(r \vee r')}\\
&\leq \frac{C''}{\lambda}
\int_0^T \d r\int_0^T \d r'
e^{-2\lambda(r \vee r')}
(C_{g,2}(r,r')-C_{g,1}(r,r'))^2
=\frac{C''}{\lambda}d_\lambda(C_{g,1},C_{g,2})^2.
\end{align*}
Thus for any $\alpha \in (0,1)$ and all sufficiently large $\lambda$,
$d_\lambda(C_{\theta,1},C_{\theta,2}) \leq
\alpha \cdot d_\lambda(C_{g,1},C_{g,2})$.

The analysis of the response $R_\theta$ is similar: Fixing $s \in [0,T]$,
there exist directional derivative processes
$D_kJ_h(\cdot,s)$
and $D_{kl}^2J_h(\cdot,s)$ over $t \in [s,T]$ for which
\begin{align*}
\lim_{\eps \to 0}
\E_{\{b^t\}}\left(\frac{J_{h+\eps k}(t,s)-J_h(t,s)}{\eps}
-D_kJ_h(t,s)\right)^2&=0,\\
\qquad \lim_{\eps \to 0}
\E_{\{b^t\}}\left(\frac{D_kJ_{h+\eps l}(t,s)-D_kJ_h(t,s)}{\eps}
-D_{kl}^2J_h(t,s)\right)^2&=0.
\end{align*}
These are the unique solutions of the equations
\begin{align*}
D_kJ_h(t,s)
    &=
    \int_s^t
    \left[
        f''(\theta_h^a)(D_k\theta_h^a)J_h(a,s)
        +
        f'(\theta_h^a)D_kJ_h(a,s)
        +
        \int_s^a R_g(a,b)D_kJ_h(b,s)\d b
    \right]\d a,\\
    D_{kl}^2J_h(t,s)
    &=
    \int_s^t
    \bigg[
        \Bigl(
            f'''(\theta_h^a)(D_k\theta_h^a)(D_l\theta_h^a)
            +
            f''(\theta_h^a)D_{kl}^2\theta_h^a
        \Bigr)J_h(a,s)
        +
        f''(\theta_h^a)(D_k\theta_h^a)D_lJ_h(a,s)
        \\
        &\qquad
        +
        f''(\theta_h^a)(D_l\theta_h^a)D_kJ_h(a,s)
                +
        f'(\theta_h^a)D_{kl}^2J_h(a,s)
        +
        \int_s^a R_g(a,b) D_{kl}^2J_h(b,s) \d b
    \bigg]\d a .
\end{align*}
Analogously to \eqref{eq:second-deriv-representation},
we may represent these solutions as
\[
    D_kJ_h(t,s)
    =
    \int_0^T M_h(t,s;r)k^r\,\d r,\qquad
    D_{kl}^2J_h(t,s)
    =
    \int_0^T\int_0^T
    L_h(t,s;r,r')k^rl^{r'}\,\d r\,\d r',
\]
where $M_h,L_h$ are the response processes
\begin{align*}
M_h(t,s;r)
    &=
    \int_s^t
    \left[
        f''(\theta_h^a)J_h(a,r)J_h(a,s)
        +
        f'(\theta_h^a)M_h(a,s;r)
        +
        \int_s^a R_g(a,b)M_h(b,s;r)\d b
    \right]\d a,\\
    L_h(t,s;r,r')
    &=
    \int_s^t
    \bigg[
        \Bigl(
            f'''(\theta_h^a)J_h(a,r)J_h(a,r')
            +
            f''(\theta_h^a)K_h(a,r,r')
        \Bigr)J_h(a,s)
        +
        f''(\theta_h^a)J_h(a,r)M_h(a,s;r')
        \\
        &\qquad
        +
        f''(\theta_h^a)J_h(a,r')M_h(a,s;r)
        +
        f'(\theta_h^a)L_h(a,s;r,r')
        +
        \int_s^a R_g(a,b)L_h(b,s;r,r')\d b
    \bigg]\d a.
\end{align*}
Applying the finite-dimensional Gaussian interpolation formula, we have
\begin{align*}
    \frac{\d}{\d u}\E J_{g_{u,N}}(t,s)
    &=
    \frac12
    \sum_{i,j=1}^N
    \bigl\langle (Q_1-Q_0)e_i,e_j\bigr\rangle_{\cH}
    \,
    \E\left[
        \partial_\eps \partial_{\eps'}J_{g_{u,N}+\eps e_i+\eps'
e_j}(t,s)\big|_{\eps=\eps'=0}
    \right]
    \\
    &=
    \frac12
    \sum_{i,j=1}^N
    \bigl\langle (Q_1-Q_0)e_i,e_j\bigr\rangle_{\cH}
    \,
    \E\left[
        \int_0^T\int_0^T
        L_{g_{u,N}}(t,s;r,r')
        e_i^re_j^{r'}
        \,\d r\,\d r'
    \right].
\end{align*}
Integrating over $u \in [0,1]$, taking the limit $N \to \infty$, and applying
that $L_{g_u}(t,s;r,r')$ is uniformly bounded and is equal to 0 unless
$r,r' \leq t$, we get
\[|(R_{\theta,2}-R_{\theta,1})(t,s)|
=\left|\lim_{N \to \infty} \E J_{g_{1,N}}(t,s)-\E J_{g_{0,N}}(t,s)\right|
\leq C \int_0^t\int_0^t
|(C_{g,2}-C_{g,1})(r,r')|\d r\,\d r'.
\]
Then, for constants $C,C'>0$ depending on $T$ but not $\lambda$,
\begin{align*}
d_\lambda(R_{\theta,1},R_{\theta,2})
&=\sup_{0 \leq s \leq t \leq T}
e^{-\lambda t}|R_{\theta,1}(t,s)-R_{\theta,2}(t,s)|\\
&\leq C\sup_{t \in [0,T]}
\int_0^t\int_0^t
e^{-\lambda t+\lambda(r \vee r')}
e^{-\lambda (r \vee r')}|C_{g,2}(r,r')-C_{g,1}(r,r')|\d r\,\d r'\\
&\leq C\sup_{t \in [0,T]}
\left(\int_0^t \int_0^t
e^{-2\lambda t+2\lambda(r \vee r')} \d r\,\d r'\right)^{1/2}
d_\lambda(C_{g,1},C_{g,2}) \leq
\frac{C'}{\sqrt{\lambda}}d_\lambda(C_{g,1},C_{g,2}).
\end{align*}
Combining these bounds for $C_\theta$ and $R_\theta$,
for any $\alpha \in (0,1)$ and all large enough $\lambda>0$,
\begin{equation}\label{eq:contractionsameRg}
d_\lambda(Y_1,Y_2) \leq \alpha \cdot d_\lambda(C_{g,1},C_{g,2})
\end{equation}

The lemma then follows from applying \eqref{eq:contractionsameCg} to bound
$d_\lambda(\cT_{g \to \theta}(C_{g,1},R_{g,1}),
\cT_{g \to \theta}(C_{g,1},R_{g,2}))$, and then
\eqref{eq:contractionsameRg} to bound
$d_\lambda(\cT_{g \to \theta}(C_{g,1},R_{g,2}),
\cT_{g \to \theta}(C_{g,2},R_{g,2}))$.
\end{proof}

\begin{lemma}[Lipschitz continuity of $\cT_{\theta \to g}$]
\label{lem:contraction-Lipschitz}
    Let $Y_i=(C_{\theta,i},R_{\theta,i}) \in \cS_\theta^+(T,D)$ for $i=1,2$, and
let $X_i=(C_{g,i},R_{g,i})=\cT_{\theta\to g}(Y_i)$. Then there exists a constant
$C>0$ not depending on $\lambda$ such that for all sufficiently large
$\lambda>0$, 
    \[
    d_\lambda(X_1,X_2) \leq C\cdot d_\lambda(Y_1,Y_2).
    \]
\end{lemma}

\begin{proof}
We have
\[
\begin{aligned}
  C_{g,1}(t,s)-C_{g,2}(t,s)
  &=
  \sum_{p,q\ge 0}
  \kappa_{p+q+2}
  \bigl(
    R_{\theta,1}^{\ast p}\ast (C_{\theta,1}-C_{\theta,2})
    \ast \bar R_{\theta,1}^{\ast q}
  \bigr)(t,s) \\
  &\hspace{0.5in}
  +
  \sum_{p,q\ge 0}
  \kappa_{p+q+2}
  \bigl[
    \bigl(R_{\theta,1}^{\ast p}\ast C_{\theta,2}\ast \bar R_{\theta,1}^{\ast q}\bigr)(t,s) -
    \bigl(R_{\theta,2}^{\ast p}\ast C_{\theta,2}\ast \bar R_{\theta,2}^{\ast q}\bigr)(t,s)
  \bigr].
\end{aligned}
\]
Consequently, 
\(
  d_\lambda(C_{g,1},C_{g,2}) \leq I_C+I_R,
\)
where
\begin{align*}
  I_C
  &:=\sum_{p,q\ge 0}
  |\kappa_{p+q+2}| \cdot
  \big\|
    R_{\theta,1}^{\ast p}\ast (C_{\theta,1}-C_{\theta,2})
    \ast \bar R_{\theta,1}^{\ast q}\big\|_\lambda,\\
  I_R
  &:=\sum_{p,q\ge 0}
  |\kappa_{p+q+2}| \cdot
  \big\|\bigl(R_{\theta,1}^{\ast p}\ast C_{\theta,2}\ast \bar R_{\theta,1}^{\ast
q}\bigr)(t,s)-\bigl(R_{\theta,2}^{\ast p}\ast C_{\theta,2}\ast \bar
R_{\theta,2}^{\ast q}\bigr)\big\|_\lambda.
\end{align*}
We first bound \(I_C\). Note that
$|R_{\theta,1}(t,s)|\le \Phi_{R_\theta}(t-s)$ for all $0\le s\le t\le T$, and
\[\|\Phi_{R_\theta}^{\ast p}\|_{L^1([0,T])}
\le \frac{(\Phi_{R_\theta}(T) T)^{p}}{p!}.\]
For $p,q \geq 1$,
\begin{align*}
&\big\|
    R_{\theta,1}^{\ast p}\ast (C_{\theta,1}-C_{\theta,2})
    \ast \bar R_{\theta,1}^{\ast q}\big\|_\lambda^2\\
    &=\int_0^T \int_0^T e^{-2\lambda(t \vee s)}
  \left(\int_0^t \d t'\int_0^s \d s'
  R_{\theta,1}^{*p}(t,t')
  (C_{\theta,1}-C_{\theta,2})(t',s')
  R_{\theta,1}^{*q}(s,s')\right)^2\d t\, \d s\\
  &\leq C \int_0^T \int_0^T e^{-2\lambda(t \vee s)}
  \int_0^t \d t'\int_0^s \d s'
  R_{\theta,1}^{*p}(t,t')^2
  (C_{\theta,1}-C_{\theta,2})(t',s')^2
  R_{\theta,1}^{*q}(s,s')^2\d t\, \d s\\
  &\leq C\int_0^T \d t'\int_0^T \d s'
  e^{-2\lambda(t' \vee s')}
  (C_{\theta,1}-C_{\theta,2})(t',s')^2
  \int_{t'}^T  \d t \int_{s'}^T \d s
  R_{\theta,1}^{*p}(t,t')^2
  R_{\theta,1}^{*q}(s,s')^2\\
  &\leq C'\|C_{\theta,1}-C_{\theta,2}\|_\lambda^2 \|\Phi_{R_\theta}^{*p}\|_{L^1([0,T])}^2
  \|\Phi_{R_\theta}^{*q}\|_{L^1([0,T])}^2,
\end{align*}
where $C,C'>0$ are $T$-dependent constants independent of $p,q$. The same bound holds for $p=0$ and/or $q=0$. Thus
\begin{align*}
  I_C 
  &\le C\sum_{p,q\ge 0} |\kappa_{p+q+2}| \frac{(\Phi_{R_\theta}(T) T)^p}{p!} \frac{(\Phi_{R_\theta}(T) T)^q}{q!} \|C_{\theta,1}-C_{\theta,2}\|_\lambda
\leq C'd_\lambda(C_{\theta,1},C_{\theta,2})
\end{align*}
for constant $C,C'>0$ not depending on $\lambda$.

We now bound \(I_R\). First,
we have $|C_{\theta,2}(t,s)|
  \le C_{\theta,2}(t,t)^{1/2} C_{\theta,2}(s,s)^{1/2} \le 
  \Phi_{C_\theta}(T)$.
Next, for the product of response kernels, writing $t_0=t$ and $s_0=s$,
\begin{align*}
&\left|
  \prod_{a=0}^{p-1}R_{\theta,1}(t_a,t_{a+1})
  \prod_{b=0}^{q-1}R_{\theta,1}(s_b,s_{b+1})
  -
  \prod_{a=0}^{p-1}R_{\theta,2}(t_a,t_{a+1})
  \prod_{b=0}^{q-1}R_{\theta,2}(s_b,s_{b+1})
\right|\\
&\le e^{\lambda(t\vee s)}
\left(
    \sup_{0\le r\le q\le T}
    e^{-\lambda q}|R_{\theta,1}(q,r)-R_{\theta,2}(q,r)|
  \right)
\sum_{\ell=1}^{p+q}
\prod_{m\ne \ell} \Phi_{R_\theta}(\xi_m),
\end{align*}
where \(\{\xi_m\}_{m=1}^{p+q}\) denotes the collection of increments
  $t_0-t_1,\ldots,t_{p-1}-t_p,
  s_0-s_1,\ldots,s_{q-1}-s_q$.
(The empty product is interpreted as \(1\).) 
Hence, integrating over $t_1,\ldots,t_p$ and $s_1,\ldots,s_q$,
\[
  I_R
  \le
  Cd_\lambda(R_{\theta,1},R_{\theta,2})
  \int_0^T\int_0^T
    \sum_{p,q\ge 0}
    |\kappa_{p+q+2}|(p+q)\Phi_{R_\theta}(T)^{p+q-1}
    \frac{t^p}{p!}\frac{s^q}{q!}\,
  \d t \d s.
\]
This shows
$I_R \le C d_\lambda(R_{\theta,1}, R_{\theta,2})$ for a constant $C>0$ not
depending on $\lambda$. Therefore,
\begin{align*}
   d_{\lambda}(C_{g,1},C_{g,2}) \le C d_\lambda(Y_1,Y_2).
  \end{align*}
  By \eqref{eq:dmft-Rg} and a similar argument omitted for brevity,
  $d_{\lambda}(R_{g,1}, R_{g,2}) \le Cd_{\lambda}(R_{\theta,1}, R_{\theta,2})$.
Thus
    $d_\lambda(X_1,X_2) \leq Cd_\lambda(Y_1,Y_2)$.
\end{proof}

Fixing $T>0$, define the projection
\[\cP_D:\cS_g(T,D) \to \cS_g^+(T,D),
\qquad (C_g,R_g) \mapsto (C_g^+,R_g)\]
that preserves $R_g$ and projects $C_g$ onto the subset $\cS_g^+(T,D)$ of the positive-semidefinite cone
\begin{equation}\label{eq:hilbertprojection}
C_g^+=\operatorname{arg\,min}_{C \in \cS_g^+(T,D)} \|C-C_g\|_\lambda
\end{equation}
in the weighted $L^2$-metric $d_\lambda$. We note that if $(C_g^{(m)},R_g^{(m)}) \in \cS_g^+(T,D)$
for $m=1,2,\ldots$ is a Cauchy sequence under $d_\lambda$,
then by completeness of the weighted $L^\infty$-metric
$d_\lambda(R_{g,1},R_{g,2})$, there exists $R_g:[0,T] \times [0,T] \to \R$ for
which
\begin{equation}\label{eq:Rgcauchy}
\lim_{m \to \infty} \sup_{s,t \in [0,T]} |R_g^{(m)}(t,s)-R_g(t,s)|=0.
\end{equation}
By completeness of the weighted $L^2$-metric $d_\lambda(C_{g,1},C_{g,2})$,
there also exists $C_g:[0,T] \times [0,T] \to \R$ for which
$\|C_g\|_\lambda<\infty$
and $\lim_{m \to \infty} d_\lambda(C_g^{(m)},C_g)=0$. 
Then the continuity condition \eqref{eq:Cg-cont2} for each kernel $C_g^{(m)}$
implies there exists
a modification $\tilde C_g:[0,T] \times [0,T] \to \R$ of $C_g$ (equal to $C_g$
Lebesgue-a.e.) for which
\begin{equation}\label{eq:Cgcauchy}
\lim_{m \to \infty} \sup_{s,t \in [0,T]} |C_g^{(m)}(t,s)-\tilde C_g(t,s)|=0.
\end{equation}
The statements \eqref{eq:Rgcauchy} and \eqref{eq:Cgcauchy} imply that $R_g$ and $\tilde C_g$ satisfy all conditions of
$\cS_g^+(T,D)$, so $(\tilde C_g,R_g) \in \cS_g^+(T,D)$,
and $d_\lambda((C_g^{(m)},R_g^{(m)}),(\tilde C_g,R_g)) \to 0$.
Then $\cS_g^+(T,D)$ is a
closed, convex subset of $\cS_g(T,D)$ under $d_\lambda$, so the
projection \eqref{eq:hilbertprojection} is well-defined. Furthermore, by the
non-expansive property of projections in a Hilbert space,
\begin{equation}\label{eq:projectioncontractive}
d_\lambda(C_{g,1}^+,C_{g,2}^+) \leq d_\lambda(C_{g,1},C_{g,2}).
\end{equation}

Lemma \ref{lemma:well_defined} then implies that
\[\cT_{g \to g}^D:=\cP_D \circ \cT_{\theta \to g} \circ \cT_{g \to \theta}\]
is a well-defined mapping from $\cS_g^+(T,D)$ to itself. The following corollary
collects the preceding results to show that this mapping is contractive in
$d_\lambda$ and has a unique fixed point in $\cS_g^+(T,D)$.
 
\begin{corollary}\label{cor:continuous-dmft-existunique-with-projection}
Fix any $T>0$ and finite subset $D \subset [0,T]$. Then
$\cT_{g \to g}^D:=\cP_D \circ \cT_{\theta \to g} \circ \cT_{g \to \theta}$
has a unique fixed point in $\cS_g^+(T,D)$.
\end{corollary}
\begin{proof}
The preceding argument verifies that 
$\cS_g^+(T,D)$ is complete under the metric $d_\lambda$. Lemmas
\ref{lem:contraction-Rtheta} and \ref{lem:contraction-Lipschitz}, together with
\eqref{eq:projectioncontractive}, show that
$\cT_{g\to g}^D$ is a contraction on $\cS_g^+(T,D)$ with respect to $d_\lambda$,
for any sufficiently large $\lambda>0$. Thus it has a unique fixed point by the
Banach fixed point theorem.
\end{proof} 

\subsection{Discrete-time dynamical mean-field limit}

We next prove a dynamical mean-field limit for a discretized version of the
dynamics, and identify continuous-time embeddings of its correlation and response
kernels as fixed points of a discretized version of the preceding mappings.

Fix a discretization step size $\eta>0$. Given $\X$, $\btheta^0$, and the
Brownian motion $\{\b^t\}_{t \geq 0}$ defining \eqref{eq:dynamics},
set $\b_\eta^k=\b^{k\eta}$ and define the discretized dynamics 
\begin{equation}\label{eq:disc-Langevin}
\btheta^{k+1}_\eta-\btheta^{k}_\eta=\eta[f(\btheta^{k}_\eta)+\X\btheta^{k}_\eta]
+ \sqrt{2\gamma} (\b^{k+1}_\eta - \b^{k}_\eta),
\qquad \btheta_\eta^0=\btheta^0.
\end{equation}
For its mean-field limit, let $\theta^0$ and $\{b^t\}_{t \geq 0}$ be as in
Definition \ref{def:DMFT}, and set
\[\theta_\eta^0=\theta^0, \qquad b_\eta^k=b^{k\eta}.\]
Given the processes $\{\theta_\eta^i\}_{0 \leq i \leq k}$
and $\{\frac{\partial \theta_\eta^i}{\partial g_\eta^j}\}_{0 \leq j<i \leq k}$
up to step $k$, define correlation and response matrices 
$\C_\theta,\bR_\theta,\C_g,\bR_g \in \R^{(k+1) \times (k+1)}$ by
\begin{equation}\label{eq:discreteCR}
\begin{gathered}
\C_\theta=\Big(\E[\theta^{i}_\eta\theta^{j}_\eta]\Big)_{0 \leq i,j \leq k},
\qquad \bR_\theta=\Big(\1_{i>j}\E\Big[\frac{\partial
\theta^i_\eta}{\partial g^j_\eta}\Big]\Big)_{0 \leq i,j \leq k},\\
\C_g=\sum_{p,q\ge 0} \kappa_{p+q+2}
\bR_\theta^p\C_\theta\bR_\theta^{q\top},
\qquad \bR_g=\sum_{p \geq 1} \kappa_{p+1}\bR_\theta^p.
\end{gathered}
\end{equation}
Here, the matrices $\bR_\theta,\bR_g$ are strictly lower-triangular.
We will verify in Theorem \ref{thm:discreteDMFT} below that
$\C_g$ is positive-semidefinite. Then, letting
\begin{equation}\label{eq:discreteg}
(g^0_\eta,\ldots,g^k_\eta) \sim \cN(0,\C_g)
\end{equation}
be independent of $\theta^0$ and $\{b^t\}_{t \geq 0}$, and writing
$\bR_g=(r_{ij})_{0 \leq i,j \leq k}$ for the entries of $\bR_g$, define
the processes up to step $k+1$ by
\begin{align}
\theta^{k+1}_\eta
&=\theta^{k}_\eta + \eta\Big[f(\theta^{k}_\eta)+\sum_{q=0}^k r_{kq}\theta^{q}_\eta+g^{k}_\eta\Big] + \sqrt{2\gamma} (b^{k+1}_\eta - b^{k}_\eta),\label{eq:DMFTdiscrete2}\\ 
\frac{\partial \theta^{k+1}_\eta}{\partial g^{j}_\eta}  
&=\frac{\partial \theta^{k}_\eta}{\partial g^{j}_\eta}+\eta\Big[f'(\theta^{k}_\eta)\frac{\partial \theta^{k}_\eta}{\partial g^{j}_\eta}
+ \sum_{q=j+1}^k r_{kq}\frac{\partial \theta^{q}_\eta}{\partial
g^{j}_\eta}\Big] \text{ for } 0 \leq j < k,
\qquad \frac{\partial \theta^{k+1}_\eta}{\partial
g^{k}_\eta}=\eta.\label{eq:DMFTdiscreteresponse}
\end{align}
These definitions (\ref{eq:discreteCR}--\ref{eq:DMFTdiscreteresponse})
are understood iteratively in time for $k=0,1,2,\ldots$, and
\eqref{eq:DMFTdiscreteresponse} is the usual derivative
of \eqref{eq:DMFTdiscrete2} in $g_\eta^j$ (by the chain rule). 
Note that $\C_\theta,\bR_\theta,\C_g,\bR_g$ up to each step $k \geq 0$
are the upper-left corners of those up to step $k+1$; thus we may write
unambiguously $(\C_\theta)_{ij},(\bR_\theta)_{ij}$ etc.\ for their entries.

We denote also
\begin{equation}\label{eq:discretegk}
\g^{k}_\eta = \X \btheta^{k}_\eta - \sum_{q=0}^k r_{kq} \btheta^{q}_\eta.
\end{equation}

\begin{theorem}\label{thm:discreteDMFT}
The matrix $\C_g$ is positive-semidefinite. 
For any fixed integer $K \ge 0$, almost surely as $n \to \infty$,
\begin{align}\label{eq:AMP-SE}
&\frac{1}{n}\sum_{i=1}^n
\delta_{(\theta^{0}_{\eta,i},\ldots,\theta^{K}_{\eta,i},g^{0}_{\eta,i},\ldots,g^{K}_{\eta,i},b^{0}_{\eta,i},\ldots,b^{K}_{\eta,i})}
\to \Law(\theta^{0}_\eta,\ldots,\theta^{K}_\eta, g^{0}_\eta,\ldots,g^{K}_\eta,
b^{0}_\eta,\ldots,b^{K}_\eta)
\end{align}
weakly and in Wasserstein-2, where the right side denotes the joint law under
(\ref{eq:discreteg}--\ref{eq:DMFTdiscrete2}). Furthermore,
for any fixed $0 \leq j \leq k$, almost surely
\begin{equation}\label{eq:discretekernelconvergence}
\frac{1}{n}\langle \btheta_\eta^{k\top}\btheta_\eta^j \rangle
\to (\C_\theta)_{k,j},
\qquad \frac{1}{n}\langle \g_\eta^{k\top}\g_\eta^j \rangle
\to (\C_g)_{k,j},
\qquad
\frac{1}{n}\langle \btheta_\eta^{k\top}
\b_\eta^j \rangle \to \sqrt{2\gamma}\sum_{q=0}^{j-1}(\bR_\theta)_{k,q}.
\end{equation}
\end{theorem}

\begin{proof}
Set $\v^0=\btheta^0$ and $\v^k=\btheta^{k}_\eta-\btheta^{k-1}_\eta$, so that
$\btheta^{k}_\eta=\v^0+\ldots+\v^k$. 
It is readily checked that for any deterministic scalar
values $\{a_{jk}\}_{0 \leq j \leq k}$,
the dynamics \eqref{eq:disc-Langevin} may be equivalently written as
\begin{equation}\label{eq:discreteAMP}
\begin{aligned}
\z^k&=\X\v^k-\sum_{j=0}^k a_{jk}\v^j,\\
\v^{k+1}&=\underbrace{\eta\left[f\Big(\sum_{j=0}^k \v^j\Big)+\sum_{j=0}^k
\Big(\sum_{i=0}^j a_{ij}\v^i+\z^j\Big)\right] +
\sqrt{2\gamma}(\b_\eta^{k+1}-\b_\eta^k)}_{:=v_{k+1}(\z^0,\ldots,\z^k,\v^0,\b_\eta^1,\ldots,\b_\eta^{k+1})},
\end{aligned}
\end{equation}
where $v_{k+1}:\R^{2k+3} \to \R$ is a scalar function defined recursively for
$k=0,1,2,\ldots$ so that
$\v^{k+1}=v_{k+1}(\z^0,\ldots,\z^k,\v^0,\b_\eta^1,\ldots,\b_\eta^{k+1})$ holds.
Since $f:\R \to \R$ is Lipschitz and continuously-differentiable,
each function $v_{k+1}$ is also Lipschitz and continuously-differentiable
in all arguments.

We define iteratively in time a corresponding AMP state evolution:
Let $v^0=\theta^0$. Given the law of Gaussian variables $(z^0,\ldots,z^{k-1})$
independent of $\theta^0$ and $\{b^t\}_{t \geq 0}$, define
\begin{align*}
\bPhi
&=\Big(\1_{i<j}\E\Big[
\frac{\partial}{\partial z^i}
v_j(z^0,\ldots,z^{j-1},v^0,b_\eta^1,\ldots,b_\eta^j)
\Big]
\Big)_{0 \leq i,j \leq k},\\
\bDelta
&=\Big(
\E[
v_i(z^0,\ldots,z^{i-1},v^0,b_\eta^1,\ldots,b_\eta^i)
v_j(z^0,\ldots,z^{j-1},v^0,b_\eta^1,\ldots,b_\eta^j)]
\Big)_{0 \leq i,j \leq k}
\end{align*}
where $v_i(z^0,\ldots,z^{i-1},v^0,b_\eta^1,\ldots,b_\eta^i) \equiv v^0$ for
$i=0$. Then define the AMP Onsager matrix and state evolution covariance by
\[
\A=\sum_{p \geq 1} \kappa_{p+1}\bPhi^p,
\qquad
\bSigma=\sum_{p,q \geq 0} \kappa_{p+q+2}
\bPhi^{p\top}\bDelta\bPhi^q,
\]
let $(z^0,\ldots,z^k) \sim \cN(0,\bSigma)$ (independent of $\theta^0$ and
$\{b^t\}_{t \geq 0}$), and let
\begin{equation}\label{eq:vkdef}
v^k=v_k(z^0,\ldots,z^{k-1},v^0,b_\eta^1,\ldots,b_\eta^k).
\end{equation}
Note that $\A$ is strictly upper-triangular, and the matrices
$\bPhi,\bDelta,\A,\bSigma$ up to step $k$ are the upper-left corners of those up
to step $k+1$.
Writing $\A=(a_{ij})_{0 \leq i,j \leq k}$ and choosing the
coefficients $a_{ij}$ in \eqref{eq:discreteAMP} to be the entries of $\A$,
\cite[Theorem 4.3]{fan2022approximate} then implies that $\bSigma$
above is positive-semidefinite (hence the above law of $(z^0,\ldots,z^k)$
is well-defined), and for any fixed $K$, almost surely in Wasserstein-2,
\begin{align}\label{eq:AMPconvergence}
\lim_{n \to \infty}
\frac{1}{n}\sum_{i=1}^n
\delta_{(v_i^0,\ldots,v_i^K,z_i^0,\ldots,z_i^K,b_{\eta,i}^0,\ldots,b_{\eta,i}^K)}
\to 
\Law\left(v^0,\ldots,v^K,z^0,\ldots,z^K,b_\eta^0,\ldots,b_\eta^K\right).
\end{align}

To check that this implies \eqref{eq:AMP-SE},
recall $\btheta_\eta^k=\v^0+\ldots+\v^k$, define
$\g_\eta^k=\z^0+\ldots+\z^k$, and define also
from the above AMP state evolution variables
\begin{align*}
\theta^{j}_\eta:=v^0+\ldots+v^j,
\qquad g^{j}_\eta:=z^0+\cdots+z^j.
\end{align*}
Denote the upper-triangular
matrix $\bE=(\1_{i\leq j})_{0\leq i,j\leq k}$. Comparing the above definitions
of $\bPhi,\bDelta,\A,\bSigma$ with the
definitions (\ref{eq:discreteCR}), we then have
\[\bR_\theta^\top=\bE^{-1}\bPhi\bE,
\qquad \C_\theta=\bE^\top\bDelta\bE,
\qquad \bR_g^\top=\bE^{-1}\A\bE,
\qquad \C_g=\bE^\top\bSigma\bE\]
and $(g_\eta^0,\ldots,g_\eta^k) \sim \cN(0,\C_g)$. In particular, $\C_g$ is
positive-semidefinite as claimed. Noting that $\A\bE=\bE\bR_g^\top$
and writing $\bR_g=(r_{ij})_{0 \leq i,j \leq k}$,
the recursions for $\z^k$ and $\v^{k+1}$ in \eqref{eq:discreteAMP} then imply
\begin{align*}
\g_\eta^k&=\X(\v^0+\ldots+\v^k)-\sum_{j=0}^k r_{kj}(\v^0+\ldots+\v^j)\\
&=\X\btheta_\eta^k-\sum_{j=0}^k r_{kj}\btheta_\eta^j,\\
\btheta_\eta^{k+1}-\btheta_\eta^k
=\v^{k+1}&=\eta\Big[
f(\v^0+\cdots+\v^k)
+\sum_{j=0}^k r_{kj}(\v^0+\cdots+\v^j)
+\g_\eta^k \Big] + \sqrt{2\gamma}(\b_\eta^{k+1}-\b_\eta^k)\\
&=\eta\Big[f(\btheta_\eta^k)+\sum_{j=0}^k r_{kj}\btheta_\eta^j+\g_\eta^k
\Big] + \sqrt{2\gamma}(\b_\eta^{k+1}-\b_\eta^k).
\end{align*}
Thus $\g_\eta^k$ coincides with the definition \eqref{eq:discretegk}, and we
have analogously from the definition of $v^{k+1}$ in \eqref{eq:vkdef} that
\begin{align*}
\theta_\eta^{k+1}-\theta_\eta^k
=v^{k+1}=\eta\Bigg[f(\theta_\eta^k)+\sum_{j=0}^k r_{kj}\theta_\eta^j
+g_\eta^k\Bigg] + \sqrt{2\gamma}(b_\eta^{k+1}-b_\eta^k),
\end{align*}
coinciding with \eqref{eq:DMFTdiscrete2}. Thus
$\{\theta_\eta^k,g_\eta^k,b_\eta^k\}_{k \geq 0}$ have the same joint law as defined in
(\ref{eq:discreteg}--\ref{eq:DMFTdiscrete2}), and the Wasserstein-2
convergence \eqref{eq:AMPconvergence} implies \eqref{eq:AMP-SE} as desired.

The first two statements of \eqref{eq:discretekernelconvergence} follow
immediately from \eqref{eq:AMP-SE} and dominated convergence to take
the expectation over $\{\b^t\}_{t \geq 0}$. Also from \eqref{eq:AMP-SE} and
dominated convergence, for any $j \leq k$,
\[
\frac{1}{n}\langle \btheta_\eta^{k\top}\b_\eta^j \rangle
\to \E [\theta^k_\eta b^j_\eta].\]
Denoting $w^q=b_\eta^{q+1}-b_\eta^q \sim \cN(0,\eta)$,
Gaussian integration-by-parts gives
\[
\E [\theta^k_\eta (b^{q+1}_\eta-b^q_\eta)] = \E [\theta^k_\eta w^q] = \eta\,\E
\Big[\frac{\partial \theta^k_\eta}{\partial w^q}\Big]
=\sqrt{2\gamma}\,\E\Big[\frac{\partial \theta^k_\eta}{\partial g^q_\eta}\Big]\]
where the last step holds by the form of the equation \eqref{eq:DMFTdiscrete2}.
Noting that $\E[\theta_\eta^k b_\eta^0]=0$ and summing over $q=0,\ldots,j-1$
shows the last statement of \eqref{eq:discretekernelconvergence}.
\end{proof}

%An important fact that deserves attention 
%from the above proof is that \begin{align}
%  \C_k &=\bE_k^\top\bSigma_k\bE_k
%=\sum_{p,q\ge 0} \kappa_{p+q+2}
%[(\bE_k^{-1}\bPhi_k\bE_k)^p]^\top(\bE_k^\top\bDelta_k\bE_k)
%(\bE_k^{-1}\bPhi_k\bE_k)^{q} \succeq 0,
%\end{align}
%as $\bSigma_k$ is positive semidefinite from the AMP state evolution, which is
%not a priori clear from the definition of $\C_k$ in \eqref{eq:RkCkupdate2}.
 
We next identify piecewise-constant continuous-time embeddings of the
discrete-time mean-field limit as fixed points of a discretized system of
mappings. Still fixing the discretization step size $\eta>0$,
for $t\geq 0$, define
\begin{align*}
  \lfloor t\rfloor
  = \max\{k\eta: k\in \mathbb Z,\ k\eta\leq t\},
\qquad \lceil t\rceil
  = \lfloor t\rfloor+\eta,
\qquad
  [t]
  = \lfloor t\rfloor/\eta \in \mathbb{Z}.
\end{align*}
We use the convention that $\int_s^t f(r)\d r=0$ if $t<s$.
Fixing $T>0$, define the discontinuity set
\[D_\eta=[0,T] \cap \eta\mathbb{Z}.\]
Let $\cS_\theta^{\eta,+} \subset \cS_\theta^+(T,D_\eta)$ and $\cS_g^\eta \subset
\cS_g(T,D_\eta)$ be the subsets of the spaces
$\cS_\theta^+(T,D_\eta),\cS_g(T,D_\eta)$ defined in Section \ref{subsec:spaces},
for which the kernels have the piecewise-constant structure
\begin{equation}\label{eq:piecewiseconstant}
C(t,s)=C(\lfloor t \rfloor,\lfloor s \rfloor),
\qquad R(t,s)=\1_{t \geq s}\,R(\lfloor t \rfloor,\lfloor s \rfloor)
\end{equation}
for $(C,R)=(C_\theta,R_\theta)$ or $(C,R)=(C_g,R_g)$, respectively. We define a map
\begin{align}\label{eq:Tthetageta}
\cT_{\theta \to g}^\eta:\cS_\theta^{\eta,+} \to \cS_g^\eta,
\qquad (C_\theta,R_\theta) \mapsto (C_g,R_g)
\end{align}
by
\begin{align*}
R_g(t,s)&=\kappa_2R_\theta(t,s)
+\sum_{p \geq 2}\kappa_{p+1}\int_{\lceil s \rceil}^{\lfloor t \rfloor} \d t_1
\int_{\lceil s \rceil}^{\lfloor t_1 \rfloor} \d t_2
\ldots \int_{\lceil s \rceil}^{\lfloor t_{p-2} \rfloor} \d t_{p-1}\\
&\hspace{3in}R_\theta(t,t_1)R_\theta(t_1,t_2)\ldots R_\theta(t_{p-1},s),\\
C_g(t,s)&=\sum_{p,q \geq 0} \kappa_{p+q+2}\int_0^{\lfloor t \rfloor} \d t_1
\ldots \int_0^{\lfloor t_{p-1} \rfloor} \d t_p
\int_0^{\lfloor s \rfloor} \d s_1
\ldots \int_0^{\lfloor s_{q-1} \rfloor} \d s_q\\
&\hspace{1in}R_\theta(t,t_1)\ldots R_\theta(t_{p-1},t_p)
R_\theta(s,s_1)\ldots R_\theta(s_{q-1},s_q) C_\theta(t_p,s_q),
\end{align*}
where $t_0=t$ and $s_0=s$. Let
\[\cS_g^{\eta,+}=\big\{\cP_{D_\eta}(C_g,R_g):(C_g,R_g) \in
\cS_g^\eta\big\} \subset \cS_g^+(T,D_\eta)\]
be the image of $\cS_g^\eta$ under the preceding projection $\cP_{D_\eta}$.
We denote
\[\cT_{\theta \to g}^{\eta,+}:\cS_\theta^{\eta,+} \to \cS_g^{\eta,+},
\qquad \cT_{\theta \to g}^{\eta,+}=\cP_{D_\eta} \circ \cT_{\theta \to g}^\eta.\]

Next, we define a map
\begin{align}\label{eq:Tgthetaeta}
  \cT_{g \to \theta}^\eta:\cS_g^{\eta,+} \to \cS_\theta^{\eta,+},
\qquad (C_g,R_g) \mapsto (C_\theta,R_\theta)
\end{align}
by
\[C_\theta(t,s)=\E[\bar\theta^t\bar\theta^s],
\qquad R_\theta(t,s)=\E\left[\frac{\partial \bar\theta^t}{\partial g^s}\right]\]
where $\{\bar\theta^t\}_{t \geq 0}$ and $\{\frac{\partial \bar\theta^t}{\partial
g^s}\}_{t \geq s \geq 0}$ solve
\begin{align}
\bar\theta^t&=\bar\theta^0+\int_0^{\lfloor t\rfloor}
  \bigg[f(\bar\theta^r)
    +\int_0^{\lfloor r\rfloor} R_g(r,q)
    \bar\theta^q\,\d q+g^{\lfloor r \rfloor}\bigg]\d r+\sqrt{2\gamma}\,b^{\lfloor
t\rfloor}, \qquad \bar\theta^0=\theta^0,
\label{eq:piecewise_discrete_dmft_theta}\\
\frac{\partial \bar\theta^t}{\partial g^s}
  &=1+\int_{\lceil s\rceil}^{\lfloor t\rfloor}
  \bigg[
    f'(\bar\theta^r)
    \frac{\partial \bar\theta^r}{\partial g^s}
    +\int_{\lceil s\rceil}^{\lfloor r\rfloor}
    R_g(r,q)\frac{\partial \bar\theta^q}{\partial g^s}\,\d q
  \bigg]\d r \text{ for } t \geq s,
\quad \frac{\partial \bar\theta^t}{\partial g^s}=0 \text{ for } t<s,\label{eq:piecewise_discrete_dmft_response}
\end{align}
and $\{g^t\}_{t \in [0,T]} \sim \GP(0,C_g)$ in
\eqref{eq:piecewise_discrete_dmft_theta} which is independent of $\theta^0$ and
$\{b^t\}_{t \geq 0}$.

\begin{lemma}\label{lemma:discretemappings}
For any fixed $T>0$ and all sufficiently small $\eta>0$, the following hold:
\begin{enumerate}[(a)]
\item For any $(C_g,R_g) \in \cS_g^{\eta,+}$, we have $\cT_{g \to
\theta}^\eta(C_g,R_g) \in \cS_\theta^{\eta,+}$.
\item For any $(C_\theta,R_\theta) \in \cS_\theta^{\eta,+}$, we have $\cT_{\theta
\to g}^\eta(C_\theta,R_\theta) \in \cS_g^\eta$
and $\cT_{\theta \to g}^{\eta,+}(C_\theta,R_\theta) \in \cS_g^{\eta,+}$.
\item Let $\C_\theta,\bR_\theta,\C_g,\bR_g$ be the matrices
\eqref{eq:discreteCR} defined up to a step
$k>T/\eta$, and define their continuous-time embeddings
\[C_\theta^\eta(t,s)=(\C_\theta)_{[t],[s]},
\qquad C_g^\eta(t,s)=(\C_g)_{[t],[s]},\]
\[R_\theta^\eta(t,s)=\begin{cases}
0 & \text{ if } s>t\\
1 & \text{ if } s \leq t,\,[s]=[t] \\
\frac{1}{\eta}\,(\bR_\theta)_{[t],[s]} & \text{ if } [s]<[t]
\end{cases},
\qquad R_g^\eta(t,s)=\begin{cases}
0 & \text{ if } s>t\\
\kappa_2 & \text{ if } s \leq t,\,[s]=[t] \\
\frac{1}{\eta}\,(\bR_g)_{[t],[s]} & \text{ if } [s]<[t].
\end{cases}
\]
Then $(C_g^\eta,R_g^\eta) \in \cS_g^{\eta,+} \cap \cS_g^\eta$ and
$(C_\theta^\eta,R_\theta^\eta) \in \cS_\theta^{\eta,+}$. Furthermore,
$\cT_{g \to \theta}^\eta(C_g^\eta,R_g^\eta)=(C_\theta^\eta,R_\theta^\eta)$
and $\cT_{\theta \to g}^\eta(C_\theta^\eta,R_\theta^\eta)=
\cT_{\theta \to g}^{\eta,+}(C_\theta^\eta,R_\theta^\eta)=(C_g^\eta,R_g^\eta)$.
\end{enumerate}
\end{lemma}
\begin{proof}
The proof is similar to \cite[Lemma 4.3]{fan2025dynamicalI}, so we will omit
some details.

Similar arguments as in Lemma \ref{lemma:well_defined}
show that $(C_g,R_g) \in \cS_g^{\eta,+}$
implies $\cT_{g \to \theta}^\eta(C_g,R_g) \in \cS_\theta^\eta(T,D_\eta)$,
and $(C_\theta,R_\theta) \in \cS_\theta^{\eta,+}$ implies
$\cT_{\theta \to g}^\eta(C_\theta,R_\theta) \in \cS_g(T,D_\eta)$. 
An induction in time checks that the processes solving
\eqref{eq:piecewise_discrete_dmft_theta}
and \eqref{eq:piecewise_discrete_dmft_response} must satisfy
\begin{equation}\label{eq:thetapiecewiseconstant}
\theta^t=\theta^{\lfloor t \rfloor},
\qquad \frac{\partial \theta^t}{\partial g^s}
=\1_{t \geq s}\frac{\partial \theta^{\lfloor t \rfloor}}{\partial g^{\lfloor s
\rfloor}}.
\end{equation}
Then $\cT_{g \to \theta}^\eta(C_g,R_g)$ satisfies \eqref{eq:piecewiseconstant},
so $\cT_{g \to \theta}^\eta(C_g,R_g)
\in \cS_\theta^{\eta,+}$. If $(C_\theta,R_\theta) \in \cS_\theta^{\eta,+}$
satisfies \eqref{eq:piecewiseconstant}, then the definition of
\eqref{eq:Tthetageta} also implies that
$\cT_{\theta \to g}^\eta(C_\theta,R_\theta)$ satisfies
\eqref{eq:piecewiseconstant}, so
$\cT_{\theta \to g}^\eta(C_\theta,R_\theta) \in \cS_g^\eta$
and hence (by definition) $\cT_{\theta \to g}^{\eta,+}(C_\theta,R_\theta) \in \cS_g^{\eta,+}$.
Thus parts (a) and (b) hold.

Letting $(C_\theta^\eta,R_\theta^\eta,C_g^\eta,R_g^\eta)$ be the embeddings of part (c),
an induction in time using the piecewise-constant properties
\eqref{eq:piecewiseconstant} and
\eqref{eq:thetapiecewiseconstant} also checks that
$\cT_{g \to \theta}^\eta(C_g^\eta,R_g^\eta)=(C_\theta^\eta,R_\theta^\eta)$
and $\cT_{\theta \to g}^\eta(C_\theta^\eta,R_\theta^\eta)=(C_g^\eta,R_g^\eta)$.
Theorem \ref{thm:discreteDMFT} shows that $\C_g$ is
positive-semidefinite; hence $C_g^\eta$ is also a
positive-semidefinite kernel on $[0,T] \times [0,T]$, so
$\cP_{D_\eta}(C_g^\eta)=C_g^\eta$. Then also
$\cT_{\theta \to g}^{\eta,+}(C_\theta^\eta,R_\theta^\eta)=(C_g^\eta,R_g^\eta)$.

Finally, similar arguments as in Lemmas \ref{lem:contraction-Rtheta}
and \ref{lem:contraction-Lipschitz} show that
$\cT_{g \to \theta}^\eta \circ \cT_{\theta \to g}^{\eta,+}:\cS_\theta^{\eta,+} \to \cS_\theta^{\eta,+}$ is contractive in
the metric $d_\lambda$ for all sufficiently large $\lambda$. Then, since
$\cS_\theta^{\eta,+}$ is also closed under $d_\lambda$,
this map has a unique fixed point $(\bar C_\theta^\eta,\bar R_\theta^\eta) \in
\cS_\theta^{\eta,+}$. Letting $(\bar C_g^\eta,\bar R_g^\eta)=
\cT_{\theta \to g}^{\eta,+}(\bar C_\theta^\eta,\bar R_\theta^\eta)
\in \cS_g^{\eta,+}$, an induction in time using \eqref{eq:Tthetageta} and
\eqref{eq:Tgthetaeta} shows that we must have
$(\bar C_\theta^\eta,\bar R_\theta^\eta)=(C_\theta^\eta,R_\theta^\eta)$
and $(\bar C_g^\eta,\bar R_g^\eta)=(C_g^\eta,R_g^\eta)$. Thus 
$(C_\theta^\eta,R_\theta^\eta) \in \cS_\theta^{\eta,+}$ and
$(C_g^\eta,R_g^\eta)\in \cS_g^{\eta,+} \cap \cS_g^{\eta}$,
showing part (c).
\end{proof}

\subsection{Proof of Theorem \ref{thm:dmft-approx}}

We proceed to show Theorem \ref{thm:dmft-approx}. The following lemma bounding
the difference between the continuous-time and discretized mappings will be
used to show part (a).

\begin{lemma}\label{lem:DMFT-mapping-discretization-distance}
There exists a constant $C>0$ depending on $T$ but not on $\lambda,\eta$
such that for all sufficiently large $\lambda>0$ and
sufficiently small $\eta>0$,
\begin{enumerate}[(a)]
\item For any $(C_g,R_g) \in \cS_g^{\eta,+}$,
\[d_\lambda\big(\cT_{g\to \theta}(C_g,R_g),\cT_{g\to \theta}^\eta(C_g,R_g)\big)
\le C\sqrt{\eta}.\]
\item For any $(C_\theta,R_\theta) \in \cS_\theta^{\eta,+}$,
\[d_\lambda\big(\cT_{\theta\to g}(C_\theta,R_\theta),\cT_{\theta\to
g}^{\eta}(C_\theta,R_\theta)\big) \le C\eta.\]
\end{enumerate}
\end{lemma}
\begin{proof}
Given $(C_g,R_g) \in \cS_g^{\eta,+}$, let $(C_\theta,R_\theta)=\cT_{g \to
\theta}(C_g,R_g)$ and $(\bar C_\theta,\bar R_\theta)=\cT_{g \to
\theta}^\eta(C_g,R_g)$.
By arguments similar to \cite[Lemma 4.4]{fan2025dynamicalI}, there exists a
coupling of $(\{\theta^t\}_{t \geq 0},\{\frac{\partial \theta^t}{\partial
g^s}\}_{t \geq s \geq 0})$ solving (\ref{eq:dmft-theta}--\ref{eq:dmft-response})
and $(\{\bar\theta^t\}_{t \geq 0},\{\frac{\partial \bar\theta^t}{\partial
g^s}\}_{t \geq s \geq 0})$ solving
(\ref{eq:piecewise_discrete_dmft_theta}--\ref{eq:piecewise_discrete_dmft_response})
for which, for all large $\lambda>0$ and small $\eta>0$,
\[\sup_{t \in [0,T]} e^{-\lambda t} \sqrt{\E(\theta^t-\bar \theta^t)^2}
\leq C\sqrt{\eta},
\qquad \sup_{0 \leq s \leq t \leq T}
e^{-\lambda t}\E\left|\frac{\partial \theta^t}{\partial g^s}
-\frac{\partial \bar\theta^t}{\partial g^s}\right| \leq C\sqrt{\eta}.\]
The second inequality implies $d_\lambda(R_\theta,\bar R_\theta) \leq
C\sqrt{\eta}$, while the first inequality and $\E(\theta^t)^2 \leq \Phi_{C_\theta}(T)$
imply
\begin{align*}
\sup_{s,t \in [0,T]} e^{-\lambda (s \vee t)}|C_\theta(s,t)-\bar C_\theta(s,t)|
&=\sup_{s,t \in [0,T]} e^{-\lambda (s \vee
t)}|\E(\theta^s\theta^t)-\E(\bar\theta^s\bar \theta^t)|\\
&\leq \sup_{s,t \in [0,T]} e^{-\lambda s}\sqrt{\E(\theta^s-\bar \theta^s)^2}
\sqrt{\E(\theta^t)^2}
+e^{-\lambda t}\sqrt{\E(\theta^t-\bar \theta^t)^2}
\sqrt{\E(\bar\theta^s)^2}\\
& \leq C'\sqrt{\eta}.
\end{align*}
Thus part (a) follows.

For part (b), given $(C_\theta,R_\theta) \in \cS_\theta^{\eta,+}$, let
$(C_g,R_g)=\cT_{\theta\to g}(C_\theta,R_\theta)$
and $(\bar C_g,\bar R_g)=\cT_{\theta\to g}^\eta(C_\theta,R_\theta)$. According
to \eqref{eq:dmft-Rg} and \eqref{eq:Tthetageta}, we know that
\[R_g-\bar R_g=\sum_{p \geq 2} \kappa_{p+1}\Big[R_\theta^{\ast p}-
R_\theta^{\lfloor\ast\rfloor p}\Big]\]
where, denoting $t=s_0$ and $s=s_p$, we have
\begin{align*}
R^{\ast p}(t,s)&= \int_{D_p(t,s)} \prod_{a=0}^{p-1}R(s_a,s_{a+1})\, \d
s_1\cdots\d s_{p-1},\quad D_p(t,s)=\left\{(s_1,\ldots,s_{p-1}):s \leq s_{p-1}\leq\cdots\leq s_1\leq t
\right\},\\
R^{\lfloor\ast\rfloor p}(t,s)&=\int_{D_{p,\eta}(t,s)} \prod_{a=0}^{p-1}
R(s_a,s_{a+1})\, \d s_1\cdots\d s_{p-1},\\
&\hspace{1in}D_{p,\eta}(t,s)=
\left\{(s_1,\ldots,s_{p-1}):s_{a+1} \leq \lfloor s_a\rfloor
\text{ for } a=0,\ldots,p-1 \right\}.
\end{align*}
Since $D_{p,\eta}(t,s)\subseteq D_p(t,s)$ and
$|R_\theta(t,s)|\le \Phi_{R_\theta}(T)$, we have the entrywise bound
\[
\left|R_\theta^{\ast p}(t,s) - R_\theta^{\lfloor \ast \rfloor p}(t,s) \right|
\leq \Phi_{R_\theta}(T)^p\, \left|D_p(t,s)\setminus D_{p,\eta}(t,s)\right|.
\]
If a point belongs to $D_p(t,s)\setminus D_{p,\eta}(t,s)$, then for at least one $a\in\{0,\ldots,p-1\}$,
$s_{a+1}>\lfloor s_a\rfloor$,
and therefore
$0\leq s_a-s_{a+1}<\eta$.
Using a union bound over the $p$ possible gaps gives
\[
\left|D_p(t,s)\setminus D_{p,\eta}(t,s)\right| \leq \frac{p(t-s)^{p-2}\eta}{(p-2)!}.
\]
Therefore, 
\begin{align*}
d_\lambda(R_g, \bar R_g)
\le \sup_{0 \leq s \leq t \leq T} |R_g(t,s)-\bar R_g(t,s)|
\leq
%\sum_{p \geq 0} |\kappa_{p+1}|\, d_\lambda((\bar R_\theta^\eta)^{\ast p}, (\bar R_\theta^\eta)^{\lfloor\ast\rfloor p}) \le |\kappa_{2}|\Phi_{R_\theta}(T)\sqrt{T\eta} +
\sum_{p \geq 2} |\kappa_{p+1}|\frac{p\Phi_{R_\theta}(T)^pT^{p-2}}{(p-2)!}\eta
\le C\eta.
\end{align*}
A similar argument shows $d_\lambda(C_g, \bar C_g) \le C\eta$.
\end{proof}

\begin{proof}[Proof of Theorem \ref{thm:dmft-approx}(a)]
Fix any $\eta>0$ small enough.
Let $X^\eta \in \cS_g^+(T,D_\eta)$ be the unique fixed
point of the projected continuous-time mapping
\[\cT_{g \to g}^\eta:=\cP_{D_\eta} \circ
\cT_{\theta \to g} \circ \cT_{g \to \theta},\]
as guaranteed by Corollary
\ref{cor:continuous-dmft-existunique-with-projection}.
This definition depends on $\eta$
through $\cP_{D_\eta}$.

Let $\bar X^\eta$ be the fixed point of the projected
discrete-time mapping
\[\bar \cT_{g \to g}^\eta:=
\cT_{\theta \to g}^{\eta,+} \circ \cT_{g \to \theta}^\eta
=\cP_{D_\eta} \circ \cT_{\theta \to g}^\eta \circ
\cT_{g \to \theta}^\eta\]
that is constructed in Lemma \ref{lemma:discretemappings}.
By the triangle inequality, we have
\begin{align*}
    d_{\lambda}(X^\eta, \bar X^\eta) &= d_{\lambda}(\cT_{g\to g}^\eta(X^\eta),
\bar\cT_{g\to g}^\eta(\bar X^\eta)) \leq d_{\lambda}(\cT_{g\to g}^\eta(X^\eta),
\cT_{g\to g}^\eta(\bar X^\eta)) + d_{\lambda}(\cT_{g\to g}^\eta(\bar X^\eta),
\bar\cT^{\eta}_{g\to g}(\bar X^\eta)).
  \end{align*}
Lemmas
\ref{lem:contraction-Rtheta} and
\ref{lem:contraction-Lipschitz}
ensure for the first term that
$d_{\lambda}(\cT_{g\to g}^\eta(X^\eta),\cT_{g\to g}^\eta(\bar X^\eta))
<\frac{1}{2}d_\lambda(X^\eta,\bar X^\eta)$, for sufficiently large $\lambda>0$.
Applying contractivity of $\cP_{D_\eta}$ with respect to $d_\lambda(\cdot)$,
the triangle inequality, and Lemmas
\ref{lem:contraction-Lipschitz} and
\ref{lem:DMFT-mapping-discretization-distance}, the second term is bounded as
\begin{align}
&d_{\lambda}(\cT^{\eta}_{g\to g}(\bar X^\eta),\bar \cT_{g\to g}^\eta(\bar
X^\eta))\notag\\
&\leq d_{\lambda}(\cT_{\theta \to g} \circ \cT_{g \to \theta}(\bar X^\eta),
\cT_{\theta \to g}^\eta \circ \cT_{g \to \theta}^\eta(\bar X^\eta))\notag\\
&\leq d_{\lambda}(\cT_{\theta \to g} \circ \cT_{g \to \theta}(\bar X^\eta),
\cT_{\theta \to g} \circ \cT_{g \to \theta}^\eta(\bar X^\eta))
+d_{\lambda}(\cT_{\theta \to g} \circ \cT_{g \to \theta}^\eta(\bar X^\eta),
\cT_{\theta \to g}^\eta \circ \cT_{g \to \theta}^\eta(\bar X^\eta))\notag\\
&\leq C\sqrt{\eta}+C\eta \leq C'\sqrt{\eta}.\label{eq:TTetabound}
\end{align}
Thus, rearranging the above shows, for some $C>0$ not depending on $\eta$,
\begin{equation}\label{eq:XXetabound}
d_{\lambda}(X^\eta, \bar X^\eta) \leq C\sqrt{\eta}.
\end{equation}

 
By \eqref{eq:XXetabound} and Lemmas \ref{lem:contraction-Rtheta} and \ref{lem:contraction-Lipschitz},
\begin{align*}
    &d_{\lambda}(X^\eta,\cT_{\theta \to g} \circ \cT_{g \to \theta}(X^\eta))\\
    &\leq d_{\lambda}(X^\eta,\bar X^\eta) +
d_{\lambda}(\bar X^\eta,\cT_{\theta \to g} \circ \cT_{g \to \theta}(\bar X^\eta)) +
d_{\lambda}(\cT_{\theta \to g} \circ \cT_{g \to \theta}(\bar X^\eta), \cT_{\theta \to g} \circ \cT_{g \to
\theta}(X^\eta)) \\ 
    &\le C''\sqrt{\eta}+
d_{\lambda}(\bar X^\eta,\cT_{\theta \to g} \circ \cT_{g \to \theta}(\bar X^\eta))
  \end{align*}
Also by the identity $\bar X^\eta=\cT_{\theta \to g}^\eta \circ \cT_{g \to
\theta}^\eta(\bar X^\eta)$ (without projection) from 
Lemma \ref{lemma:discretemappings} and the bound for the second line of
\eqref{eq:TTetabound},
\[d_{\lambda}(\bar X^\eta,\cT_{\theta \to g} \circ \cT_{g \to \theta}(\bar
X^\eta))=d_{\lambda}(\cT_{\theta \to g}^\eta \circ \cT_{g \to \theta}^\eta(\bar
X^\eta),\cT_{\theta \to g} \circ \cT_{g \to \theta}(\bar X^\eta))
\leq C'\sqrt{\eta}.\]
So for some $C>0$ not depending on $\eta$,
\begin{equation}\label{eq:Xfixedpoint}
d_{\lambda}(X^\eta,\cT_{\theta \to g} \circ \cT_{g \to \theta}(X^\eta))
\leq C\sqrt{\eta}.
\end{equation}

 For any $\delta>0$, let $D_{\eta \cup \delta} = D_{\eta} \cup D_{\delta}$, and define $X^{\delta} \in \mathcal{S}^+(T,D_{\delta})$ analogously to $X^{\eta}$. Then $X^{\eta}, X^{\delta} \in \mathcal{S}_g^+(T,D_{\eta \cup \delta})$, and 
\begin{align}
   \nonumber d_{\lambda}(X^\eta,X^\delta) &\leq d_{\lambda}(X^\eta, \mathcal{T}_{\theta \to g} \circ \mathcal{T}_{g \to \theta}(X^\eta))+ d_{\lambda}(\mathcal{T}_{\theta \to g} \circ \mathcal{T}_{g \to \theta}(X^\delta),X^\delta)\\ \nonumber& \qquad \qquad \qquad+ d_{\lambda}(\mathcal{T}_{\theta \to g} \circ \mathcal{T}_{g \to \theta}(X^\eta),\mathcal{T}_{\theta \to g} \circ \mathcal{T}_{g \to \theta}(X^\delta))\\ \nonumber& \leq C(\sqrt{\eta}+\sqrt{\delta})+ \frac{1}{2}d_{\lambda}(X^\eta,X^{\delta}).
\end{align}
Consequently, $d_{\lambda}(X^{\eta},X^{\delta}) \leq C'(\sqrt{\eta}+\sqrt{\delta})$. This shows along any sequence of values $\eta>0$ converging to 0 that
$\{X^{\eta}\}$ is a Cauchy sequence under $d_\lambda$. Then by completeness of the ambient weighted $L^\infty$ and $L^2$ spaces over $[0,T] \times [0,T]$ as in \eqref{eq:Rgcauchy} and \eqref{eq:Cgcauchy}, there
exist $R_g,C_g:[0,T] \times [0,T] \to \R$ such that $X=(C_g,R_g)$ satisfies $d_{\lambda}(X,X^\eta) \leq C\sqrt{\eta}$. Together with \eqref{eq:Xfixedpoint}, we have $d_{\lambda}(X,\mathcal{T}_{\theta \to g}\circ \mathcal{T}_{g \to \theta}(X^\eta)) \leq C\sqrt{\eta}.$   By Lemma \ref{lemma:well_defined},
the image of $\cT_{\theta \to g} \circ \cT_{g \to \theta}$ is contained in
$\cS_g(T,\emptyset)$ with discontinuity set $D=\emptyset$. Then as in \eqref{eq:Cgcauchy}, there exists a continuous modification of $C_g$ so that $X=(C_g,R_g) \in \cS_g(T,\emptyset)$. Writing $X^\eta=(C_g^\eta,R_g^\eta)$, since also $\lim_{\eta \to 0} d_\lambda(X,X^\eta)=0$ and $C_g^\eta$ is positive-semidefinite, we have that $C_g$ is positive-semidefinite, so $X \in \mathcal{S}_g^+(T,\emptyset)$. 

Finally, \begin{align*}
   \nonumber &d_{\lambda}(X, \mathcal{T}_{\theta \to g}\circ \mathcal{T}_{g \to \theta}(X))
    \\&\leq d_{\lambda}(X, X^\eta)+ d_{\lambda}(X^{\eta}, \mathcal{T}_{\theta \to g}\circ \mathcal{T}_{g \to \theta}(X^\eta)) + d_{\lambda}(\mathcal{T}_{\theta \to g}\circ \mathcal{T}_{g \to \theta}(X^\eta), \mathcal{T}_{\theta \to g}\circ \mathcal{T}_{g \to \theta}(X)) 
 \\& \leq C\sqrt{\eta}.
\end{align*}
Since $\eta>0$ is arbitrary, this shows that $X \in \cS_g^+(T,\emptyset)$ is a fixed point of $\mathcal{T}_{\theta \to g}\circ \mathcal{T}_{g \to \theta}$. Uniqueness of such a fixed point in $\cS_g^+(T,\emptyset)$ is immediate as $\mathcal{T}_{\theta \to g}\circ \mathcal{T}_{g \to \theta}$ is contractive under $d_\lambda$.

This shows that the components of Theorem \ref{thm:dmft-approx}(a) are 
well-defined up to any fixed time $T>0$, where the fixed points
$(C_\theta,R_\theta)$ and $(C_g,R_g)$ are unique in
$\cS_\theta^+(T,\emptyset)$ and
$\cS_g^+(T,\emptyset)$. Then Theorem \ref{thm:dmft-approx}(a) holds upon setting
$\cS_\theta^+$ as the space of kernels
$(C_\theta,R_\theta)$ on $[0,\infty) \times [0,\infty)$ whose
restrictions to $[0,T] \times [0,T]$ belong to $\cS_\theta^+(T,\emptyset)$ for
each $T>0$, and similarly for $\cS_g^+$.
\end{proof}

In the remainder of the proof, we fix $T>0$,
$X=(C_g,R_g) \in \cS_g^+(T,\emptyset)$ as
the preceding unique fixed point of $\cT_{\theta \to g} \circ \cT_{g \to
\theta}$, and $\bar X^\eta=(C_g^\eta,R_g^\eta) \in \cS_g^+(T,D_\eta)$ as the
fixed point of $\cT_{\theta \to g}^\eta \circ \cT_{g \to
\theta}^\eta$. To show Theorem \ref{thm:dmft-approx}(b), we record the following
continuity properties of $C_g,C_g^\eta$.
 
\begin{lemma}\label{lem:discrete_dmft_approx_continue}
For all sufficiently small \(\eta>0\),
there is a constant $C>0$ depending on $T$ but not $\eta$, such that 
\begin{align}
  |C_g(t,s)-C_g(t',s')|
  &\le
  C\left(\sqrt{|t-t'|}+\sqrt{|s-s'|}\right),
  \label{eq:Cg-holder-cont}\\
  |C_g^\eta(t,s)-C_g^\eta(t',s')|
  &\le
  C\left(\sqrt{|t-t'|+\eta}+\sqrt{|s-s'|+\eta}\right)
  \label{eq:Cg-holder-discrete}
\end{align}
for all \(s,t,s',t'\in[0,T]\).
\end{lemma}
\begin{proof}
Since $(C_g,R_g) \in \cS_g^+(T,\emptyset)$, we have from \eqref{eq:Cg-cont2} that
\[
  |C_g(t,s)-C_g(t',s')|
  \leq
  C\left(\sqrt{|t-t'|}+\sqrt{|s-s'|}\right).
\]
For \eqref{eq:Cg-holder-discrete}, let $\{\bar \theta^t,\frac{\partial
\bar \theta^t}{\partial g^s}\}$ be the processes solving the discretized
equations
(\ref{eq:piecewise_discrete_dmft_theta}--\ref{eq:piecewise_discrete_dmft_response})
for the kernels $(C_g^\eta,R_g^\eta)$. By arguments analogous to those of
Lemma \ref{lemma:well_defined},
for some $C>0$ depending on $T$ but not $\eta$, we have  
\begin{align*}
|C_\theta^\eta(t,s)-C_\theta^\eta(t',s')|
&\leq C(\sqrt{|t-t'|+\eta}+\sqrt{|s-s'|+\eta}) \text{ for all } s,t \in
[0,T],\\
|R_\theta^\eta(t,s)-R_\theta^\eta(t',s)| &\le C(|t-t'|+\eta)
\text{ for all } s \leq t,t' \leq T,
\end{align*}
and hence also  
\[|C_g^\eta(t,s)-C_g^\eta(t',s')|
  \leq
  C\left(\sqrt{|t-t'|+\eta}+\sqrt{|s-s'|+\eta}\right).\]
\end{proof}

\begin{proof}[Proof of Theorem \ref{thm:dmft-approx}(b)]
Fix $\eta>0$, and consider a piecewise-constant embedding 
of the discretized Langevin process
\eqref{eq:disc-Langevin} defined as $\bar\btheta^t = \btheta_\eta^{[t]}$.
An induction in time checks that this solves
\begin{equation}\label{eq:langevin_discret}
\bar\btheta^t=\btheta^0+\int_0^{\lfloor{t}\rfloor} [f(\bar\btheta^s)
+\X\bar\btheta^s]\d s +\sqrt{2\gamma}\,\b^{\lfloor{t}\rfloor}
\end{equation}
where $\{\b^t\}_{t \geq 0}$ is the same Brownian motion as defining the
continuous process $\{\btheta^t\}_{t \geq 0}$.
Define also a piecewise constant embedding $\bar \g^t=\g_\eta^{[t]}$ of
\eqref{eq:discretegk}, and $\bar \b^t=\b_\eta^{[t]}=\b^{[t]\eta}$.
By \cite[Lemma 4.7]{fan2025dynamicalI}, for a constant $C>0$ independent of
$\eta$, a.s.\ for all large $n$,
\begin{align*}
\sup_{t\in[0,T]} \frac{1}{\sqrt{n}}\|\b^t\|_2\leq C,
\qquad \sup_{t\in[0,T]}
\frac{1}{\sqrt{n}}\big(\|{\b^t - \b^{\lfloor{t}\rfloor}}\|_2 + \|{\b^t -
\b^{\lceil{t}\rceil}}\|_2\big) \leq C\sqrt{\eta\max(\log(1/\eta),1)}.
\end{align*}
Then by a Gr\"onwall argument (see \cite[Lemma 4.8]{fan2025dynamicalI}), a.s.\ for
all large $n$,
\begin{align}
\sup_{t\in[0,T]}
\frac{1}{\sqrt{n}}\|{\btheta^t}\|_2 \leq C,\qquad \sup_{t\in[0,T]}
\frac{1}{\sqrt{n}}\|{\btheta^t - \bar\btheta^{t}}\|_2 \leq
C\sqrt{\eta\max(\log(1/\eta),1)}\label{eq:theta-discret-norm}
\end{align}
For $\bar \g^t$, we have
\begin{align*}
    \|\g^t - \bar \g^t\|_2 &= \left\|\X(\btheta^t-\bar\btheta^t)-\left(\int_0^t
R_g(t,s)\btheta^s\d s - \int_0^{\lfloor t\rfloor} R_g(t,s)\bar\btheta^s \d
s\right)\right\|_2\\
    & \le \|\X\|_\op\|\btheta^t-\bar\btheta^t\|_2 + \Phi_{R_g}(T)\int_0^{\lfloor
t\rfloor} \|\btheta^s-\bar\btheta^s\|_2 \d s +  \Phi_{R_g}(T)\int_{\lfloor
t\rfloor}^t \|\btheta^s\|_2\d s,
\end{align*}
so also a.s.\ for all large $n$,
\begin{align*}
\sup_{t\in[0,T]}
\frac{1}{\sqrt{n}}\|{\g^t - \bar\g^{t}}\|_2 \leq C'\sqrt{\eta\max(\log(1/\eta),1)}.
\end{align*}
This implies that for any fixed $m \geq 1$ and $t_1,\ldots, t_m \in [0,T]$,
\begin{align}
&\limsup_{n\rightarrow\infty} W_2\Bigg( \frac{1}{n}\sum_{j=1}^n
\delta_{\left(\theta^{t_1}_j, \ldots, \theta^{t_m}_j,g^{t_1}_j, \ldots,
g^{t_m}_j,b^{t_1}_j, \ldots, b^{t_m}_j\right)},\notag\\
&\hspace{3.4cm} \frac{1}{n}\sum_{j=1}^n
\delta_{\left(\bar\theta^{t_1}_j, \ldots,
\bar\theta_j^{t_m},\bar g^{t_1}_j, \ldots,
\bar g_j^{t_m},\bar b^{t_1}_{j}, \ldots,
\bar b_{j}^{t_m}\right)}\Bigg)<C\sqrt{\eta\max(\log(1/\eta),1)}\text{ a.s.}\label{eq:discretelangevinW2}
\end{align}

Now let $\{\theta^t,g^t,b^t\}_{t \geq 0}$ be the components of the
continuous-time dynamical mean-field limit in Definition
\ref{def:DMFT}, defined from the fixed point $X=(C_g,R_g)$.
Let $\{\bar\theta^t\}_{t \geq 0}$ be the solution of
\eqref{eq:piecewise_discrete_dmft_theta} with the same Brownian motion
$\{b^t\}_{t \geq 0}$, defined from the discrete-time fixed
point $\bar X^\eta=(C_g^\eta,R_g^\eta)$ and $\{\bar g^t\}_{t \geq 0} \sim
\GP(0,C_g^\eta)$. Let $\bar b^t=b^{\lfloor t \rfloor}$. The bound
\eqref{eq:XXetabound} implies that for any sufficiently large $\lambda>0$,
\begin{align*}
\lim_{\eta \to 0} \sup_{0\leq s \leq t \leq T} |R_g(t,s)-R_g^\eta(t,s)|
&\leq \lim_{\eta \to 0} e^{\lambda T}d_\lambda(R_g,R_g^\eta)=0,\\
\lim_{\eta \to 0} \int_0^T \d t \int_0^T \d s\,|C_g(t,s)-C_g^\eta(t,s)|^2
&\leq \lim_{\eta \to 0} e^{2\lambda T}d_\lambda(C_g,C_g^\eta)=0.
\end{align*}
The continuity properties for $C_g,C_g^\eta$ in Lemma
\ref{lem:discrete_dmft_approx_continue} then imply that also
\[\lim_{\eta \to 0} \sup_{s,t \in [0,T]} |C_g(t,s)-C_g^\eta(t,s)|=0.\]
Then (c.f.\ \cite[Lemma~D.1]{fan2025dynamicalII}) there exists a coupling of
$\{g^t\}_{t \in [0,T]}$ and $\{\bar g^t\}_{t \in [0,T]}$ for which
\[\lim_{\eta \to 0} \sup_{t \in [0,T]} \E(g^t-\bar g^t)^2=0.\]
Under this coupling of $\{g^t\}_{t \in [0,T]}$ and $\{\bar g^t\}_{t \in [0,T]}$,
a Gr\"onwall argument (see \cite[Eq.\ (100)]{fan2025dynamicalI})
shows that the resulting coupling
of $\{\theta^t\}_{t \in [0,T]}$ and $\{\bar \theta^t\}_{t \in [0,T]}$ satisfies
\[\lim_{\eta \to 0} \sup_{t \in [0,T]} \E(\theta^t-\bar\theta^t)^2=0.\]
Hence, for any fixed $m \geq 1$ and $t_1,\ldots,t_m \in [0,T]$,
\begin{align}
&\lim_{\eta \to 0} W_2\Big(
\Law(\theta^{t_1},\ldots,
\theta^{t_m}, g^{t_1},\ldots,
g^{t_m}, b^{t_1},\ldots, b^{t_m}),\notag\\
&\hspace{1in}\Law(\bar\theta^{t_1},\ldots,
\bar\theta^{t_m},\bar g^{t_1},\ldots,
\bar g^{t_m}, \bar b^{t_1},\ldots,\bar b^{t_m})\Big)=0.\label{eq:discretedmftW2}
\end{align}
Combining \eqref{eq:discretelangevinW2}, \eqref{eq:discretedmftW2}, and Theorem
\ref{thm:discreteDMFT} which shows
\[\lim_{n \to \infty} W_2\Bigg(\frac{1}{n}\sum_{j=1}^n
\delta_{\left(\bar\theta^{t_1}_j, \ldots,
\bar\theta_j^{t_m},\bar g^{t_1}_j, \ldots,
\bar g_j^{t_m},\bar b^{t_1}_{j}, \ldots,
\bar b_{j}^{t_m}\right)},\,\Law(\bar\theta^{t_1},\ldots,
\bar\theta^{t_m},\bar g^{t_1},\ldots,
\bar g^{t_m}, \bar b^{t_1},\ldots,\bar b^{t_m})\Bigg)=0 \text{ a.s.},\]
and taking $\eta \to 0$, we have that for any fixed $m \geq 1$
and $t_1,\ldots,t_m \in [0,T]$,
\begin{align}\label{eq:pointwise-W2-convergence}
\lim_{n \to \infty}
W_2\left(\frac{1}{n}\sum_{j=1}^n \delta_{\left(\theta^{t_1}_j, \ldots,
\theta^{t_m}_j,g^{t_1}_j, \ldots, g^{t_m}_j,b^{t_1}_j, \ldots,
b^{t_m}_j\right)},\,\Law(\theta^{t_1},\ldots,
\theta^{t_m}, g^{t_1},\ldots,g^{t_m}, b^{t_1},\ldots,b^{t_m})\right)=0
\text{ a.s.}
\end{align}

We may strengthen this to Wasserstein-2 convergence over $C([0,T],\R^3)$ using
a similar tightness argument as in \cite[Theorem 2.5(b)]{fan2025dynamicalI}:
Let
\[\x=\{\x^t\}_{t \in [0,T]}
=\{(\btheta^t,\g^t,\b^t)\}_{t \in [0,T]},
\qquad
x=\{x^t\}_{t \in [0,T]}=\{(\theta^t,g^t,b^t)\}_{t \in [0,T]}.\]
Denote
\[\|x\|_\infty=\sup_{t \in [0,T]} |\theta^t|+
\sup_{t \in [0,T]} |g^t|+\sup_{t \in [0,T]} |b^t|,\]
and write the coordinates of $\x$ as $x_j=(\theta_j,g_j,b_j) \in C([0,T],\R^3)$ for $j=1,\ldots,n$.
Fix any test function $F:C([0,T],\R^3) \to \R$ satisfying
\begin{equation}\label{eq:Fpseudolipschitz}
\big|F(x)-F(x')\big|
\leq C\|x-x'\|_\infty(1+\|x\|_\infty+\|x'\|_\infty).
\end{equation}
Fix any $\eta>0$, and let $\{\tilde\x^t\}_{t \in [0,T]}$ and $\{\tilde x^t\}_{t
\in [0,T]}$ be the piecewise-linear
interpolations of $\{\tilde \x^t\}_{t \in [0,T] \cap \eta \mathbb{Z}}$
and $\{\tilde x^t\}_{t \in [0,T] \cap \eta \mathbb{Z}}$, respectively, with
knots at $[0,T] \cap \eta \mathbb{Z}$. Then
\eqref{eq:pointwise-W2-convergence} implies
\begin{align}\label{eq:pointwise-W2-convergence-joint}
\lim_{n \to \infty} \frac{1}{n}\sum_{j=1}^n F(\{\tilde x_j^t\}_{t \in [0,T]})
=\E F(\{\tilde x^t\}_{t \in [0,T]}) \text{ a.s.}
\end{align}

By the property \eqref{eq:Fpseudolipschitz} and Cauchy-Schwarz, we have
\[\left|\frac{1}{n}\sum_{j=1}^n F(\{x_j^t\}_{t \in [0,T]})
-F(\{\tilde x_j^t\}_{t \in [0,T]})\right|
\leq C\left(\frac{1}{n}\sum_{j=1}^n \|x_j-\tilde x_j\|_\infty^2\right)^{1/2}
\left(1+\frac{1}{n}\sum_{j=1}^n \|x_j\|_\infty^2\right)^{1/2}.\]
By a modulus-of-continuity bound for Brownian motion,
the bound \eqref{eq:theta-discret-norm}, and a Gr\"onwall argument
(see \cite[Theorem 2.5(b)]{fan2025dynamicalI}), for some $C>0$ and
$\alpha \in (0,1/2)$, a.s.\ for all large $n$,
\[\frac{1}{n}\sum_{j=1}^n \|b_j\|_\infty^2 \leq C,
\quad \frac{1}{n}\sum_{j=1}^n \|\theta_j\|_\infty^2 \leq C,
\quad \frac{1}{n}\sum_{j=1}^n \|g_j\|_\infty^2 \leq C,\]
\[\frac{1}{n}\sum_{j=1}^n \|b_j-\tilde b_j\|_\infty^2 \leq C\eta^{2\alpha},
\quad \frac{1}{n}\sum_{j=1}^n \|\theta_j-\tilde\theta_j\|_\infty^2 \leq
C\eta^{2\alpha},
\quad \frac{1}{n}\sum_{j=1}^n \|g_j-\tilde g_j\|_\infty^2 \leq
C\eta^{2\alpha}.\]
Applying this above and taking $\eta \to 0$,
\begin{equation}\label{eq:W2tightnesslangevin}
\lim_{\eta \to 0}
\left|\frac{1}{n}\sum_{j=1}^n F(\{x_j^t\}_{t \in [0,T]})
-F(\{\tilde x_j^t\}_{t \in [0,T]})\right|=0 \text{ a.s.}
\end{equation}
For the components of the dynamical mean-field limit,
by standard Gaussian process bounds and a Gr\"onwall argument, also
\[\E\|b\|_\infty^2 \leq C,
\quad \E\|g\|_\infty^2 \leq C,
\quad \E\|\theta\|_\infty^2 \leq C,\]
\[\E\|b-\tilde b\|_\infty^2 \leq C\eta^{2\alpha},
\quad \E\|g-\tilde g\|_\infty^2 \leq C\eta^{2\alpha},
\quad \E\|\theta-\tilde\theta\|_\infty^2 \leq C\eta^{2\alpha},\]
and hence
\begin{equation}\label{eq:W2tightnessdmft}
\lim_{\eta \to 0} \big|\E F(\{x^t\}_{t \in [0,T]})
-\E F(\{\tilde x^t\}_{t \in [0,T]})\big|=0.
\end{equation}
Combining \eqref{eq:pointwise-W2-convergence-joint}, \eqref{eq:W2tightnesslangevin},
and \eqref{eq:W2tightnessdmft} and taking $\eta \to 0$
shows
\begin{align}\label{eq:pointwise-W2-convergence-2}
\frac{1}{n}\sum_{j=1}^n F(\{x_j^t\}_{t \in [0,T]})
\rightarrow \E F(\{x^t\}_{t \in [0,T]}) \text{ a.s.}
\end{align}
This holds for any function $F:C([0,T],\R^3) \to \R$
satisfying \eqref{eq:Fpseudolipschitz}, implying the claimed Wasserstein-2
convergence in Theorem \ref{thm:dmft-approx}(b).
\end{proof}

\begin{proof}[Proof of Theorem \ref{thm:dmft-approx}(c)]
Theorem \ref{thm:dmft-approx}(b) implies, almost surely as $n \to \infty$,
\[\frac{1}{n}\sum_{i=1}^n \theta_i^t\theta_i^s \to
\E[\theta^t \theta^s]=C_\theta(t,s),
\quad \frac{1}{n}\sum_{i=1}^n g_i^tg_i^s \to
\E[g^tg^s]=C_g(t,s).\]
The first two claims of Theorem \ref{thm:dmft-approx}(c) follow from this and
dominated convergence to take an expectation over $\{\b^t\}_{t \in [0,T]}$. For
the third claim, \eqref{eq:discretelangevinW2} implies
\[\lim_{\eta \to 0}
\limsup_{n \to \infty} \left|\frac{1}{n}\sum_{i=1}^n \theta_i^t b_i^s
-\bar\theta_i^t \bar b_i^s\right|=0 \text{ a.s.}\]
Theorem \ref{thm:discreteDMFT} shows
\[\frac{1}{n}\sum_{i=1}^n \bar\theta_i^t \bar b_i^s
\to \sqrt{2\gamma}\int_0^{\lfloor s \rfloor} R_\theta^\eta(t,s) \d s \text{ a.s.}\]
The bound \eqref{eq:XXetabound} and $|R_\theta^\eta(t,s)| \leq
\Phi_{R_\theta}(T)$ imply
\[\lim_{\eta \to 0}\left|\int_0^s R_\theta(t,s) \d s 
-\int_0^{\lfloor s \rfloor} R_\theta^\eta(t,s) \d s\right|=0.\]
Then the third claim of Theorem \ref{thm:dmft-approx}(c) follows from combining
these three statements, taking $\eta \to 0$, and again applying dominated
convergence to take an expectation over $\{\b^t\}_{t \in [0,T]}$.
\end{proof}
